\documentclass[11pt,reqno]{amsart}
\usepackage[T1]{fontenc}
\usepackage{lmodern}
\usepackage{amsmath,amssymb,mathtools}
\usepackage[margin=1.1in]{geometry}
\usepackage{microtype}
\usepackage{needspace}
\usepackage[colorlinks=true,linkcolor=blue,citecolor=blue,urlcolor=blue]{hyperref}
\newtheorem{theorem}{Theorem}[section]
\newtheorem{proposition}[theorem]{Proposition}
\newtheorem{lemma}[theorem]{Lemma}
\newtheorem{corollary}[theorem]{Corollary}
\theoremstyle{definition}
\newtheorem{definition}[theorem]{Definition}
\theoremstyle{remark}
\newtheorem{remark}[theorem]{Remark}
\newtheorem*{remark*}{Remark}
\newcommand{\C}{\mathbb C}
\newcommand{\N}{\mathbb N}
\newcommand{\B}{\mathcal B}
\newcommand{\A}{\mathcal A}
\newcommand{\M}{\mathcal M}
\newcommand{\R}{\mathcal R}
\newcommand{\HH}{\mathcal H}
\newcommand{\F}{\mathcal F}
\newcommand{\bF}{\mathbb F}
%%%%%% brackets for the free field
\makeatletter
\def\moverlay{\mathpalette\mov@rlay}
\def\mov@rlay#1#2{\leavevmode\vtop{%
\baselineskip\z@skip \lineskiplimit-\maxdimen
\ialign{\hfil$#1##$\hfil\cr#2\crcr}}}
\makeatother
\def\plangle{\moverlay{(\cr<}}
\def\prangle{\moverlay{)\cr>}}
\newcommand{\freefield}[1]{\C\mathopen{\plangle}#1\mathclose{\prangle}}
\newcommand{\Div}{\operatorname{Div}}
\newcommand{\rank}{\operatorname{rank}}
\newcommand{\dom}{\operatorname{Dom}}
\newcommand{\Ker}{\operatorname{Ker}}
\newcommand{\Image}{\operatorname{Im}}
\newcommand{\Hom}{\operatorname{Hom}}
\newcommand{\hotimes}{\mathbin{\widehat{\otimes}}}
\newcommand{\id}{\operatorname{id}}
\newcommand{\dist}{\operatorname{dist}}
\title[Rationality criteria]{Rationality criteria}
\date{\today}
\subjclass[2020]{46L54, 46L52, 16K40}
\keywords{Free field, affiliated operator, dual system, conjugate variable, finite-rank commutator}
\begin{document}
\begin{abstract}
Let $X:=(X_1,\ldots,X_d)$ be a tuple in a tracial $W^*$-probability space.
We establish three criteria for affiliated operators to be rational in $X$, each with several equivalent characterizations.
The first criterion is of algebraic nature and assumes only $\Delta(X)=d$ for $X$.
From it, we derive a second criterion of an analytic nature using liberation dual systems.
When $X$ consists of normal generators and admits a dual system, we obtain a third criterion that further refines the second in terms of finite-rank fraction commutators, thereby extending the rationality criteria of Duchamp and Reutenauer, Linnell, and Miyagawa for affiliated operators.
\end{abstract}
\maketitle
% Responsible statements: C_F0068 Part A, C_F0079 v4, C_F0063, and C_F0024 with C_F0068 Part C.
\section{Introduction}\label{sec:introduction}

Evaluating a noncommutative rational function at a tuple $X$ of bounded operators can yield an unbounded affiliated operator.
The free-field realization theorem of Mai, Speicher and Yin \cite[Theorem 1.1]{MSY23} shows that evaluation at $X$ embeds the free field into the $*$-algebra of affiliated operators when $\Delta(X)=d$.
Under this condition, we first establish algebraic characterizations of rationality in $X$.
We then derive two analytic criteria: one using liberation dual systems and the other using dual systems.

Our question is motivated by Connes's rationality conjecture for free group algebras \cite[Section IV.5, Remark 3]{Con94}, a noncommutative analogue of Kronecker's criterion \cite{Kro81}.
In its operator formulation, this criterion characterizes rationality by a finite-rank commutator condition \cite[Section IV.3, Proposition 1]{Con94}.
For a free group, Connes observed that commutators with the operator arising from its Cayley tree have finite rank on the rational closure of the group algebra in the reduced group $C^*$-algebra,
and conjectured the converse.

Duchamp and Reutenauer \cite[Theorem 7]{DR97} proved this conjecture for a free group $\bF_d$ with $d$ generators ($d\geq1$),
and showed that the rational closure coincides with the division closure.
Their proof uses noncommutative rational series.
Using the bounded dual system for free Haar unitaries (see, for example, \cite[Proposition 6.3]{MSY23}),
one can express the Cayley-tree finite-rank condition as finite rank of the corresponding commutators.

Linnell \cite[Theorem 1.1]{Lin00} extends the free-group criterion to affiliated operators using left and right fraction representations.
Theorem~\ref{thm:analytic} extends this criterion to normal generators admitting a dual system in the sense of Definition~\ref{def:dual-system}, including free Haar unitaries with their bounded dual systems.
Linnell also identifies the bounded part of the affiliated division closure with the division closure in the reduced group $C^*$-algebra.
Together with this identity, Theorem~\ref{thm:analytic} recovers the criterion of Duchamp and Reutenauer.

Miyagawa \cite[Theorems 3.1 and 4.5]{Miy22} proves a semicircular analogue of these criteria, characterizing rationality by finite-rank commutators with the right annihilation operators (whose negatives form a bounded dual system in our convention).
His proof adapts the noncommutative rational-series approach of Duchamp and Reutenauer to the semicircular setting.
He also extends this characterization to affiliated operators using left and right fraction representations.

\subsection{A rationality criterion under maximal \texorpdfstring{$\Delta$}{Delta}}
Our first main result gives several equivalent characterizations under maximal $\Delta$ alone.
The quantity $\Delta$ was introduced by Connes and Shlyakhtenko in their paper \cite{CS05} on the homology of von Neumann algebras.

Throughout the paper, $(\M,\tau)$ is a tracial $W^*$-probability space, that is, a von Neumann algebra $\M$ equipped with a faithful normal tracial state $\tau$.
Let $X=(X_1,\ldots,X_d)\in\M^d$, with $d\geq1$.
Put $H:=L^2(\M,\tau)$ and $\Omega:=1$.
Inner products on $H$ are linear in the second variable, and we identify elements of $\M$ with left multiplication on $H$.
The condition $\Delta(X)=d$ means that there are no nonzero finite-rank operators $T_i\in\F(H)$ such that
\[\sum_i[\rho(X_i),T_i]=0;\]
see Section~\ref{subsec:delta-tensor-maps} for details.
Here $\rho(X_i)$ denotes the right multiplication operator by $X_i$ on $H$.

We write $L^0(\M)$ for the $*$-algebra of closed densely defined operators affiliated with $\M$.
We identify $H$ with the square-integrable affiliated operators in $L^0(\M)$ \cite[Theorem~7.3.2]{AP18}.
In particular, $H\otimes_{\C}H\subseteq L^0(\M)\otimes_{\C}L^0(\M)^{\mathrm{op}}$.
Write $\A_X:=\C\langle X_1,\ldots,X_d\rangle$ for the algebra generated by $X_1,\ldots,X_d$.
We denote its division closure in $L^0(\M)$ by $\Div_{L^0(\M)}\A_X$ and call an affiliated operator \emph{rational} when it belongs to this closure; see Definition~\ref{def:rationality}.

\begin{definition}\label{def:fraction-pair}
For $u\in L^0(\M)$, a \emph{fraction pair} is a quadruple $Q=(s,a,b,t)\in\M^4$ satisfying
\begin{equation}
 u=s^{-1}a=bt^{-1},\qquad s,t\text{ invertible in }L^0(\M).
 \label{eq:fractions}
\end{equation}
\end{definition}

\begin{definition}\label{def:fraction-commutator}
Let $D$ be a linear operator on $H$ with domain $\dom D$.
For a fraction pair $Q=(s,a,b,t)$ of an element $u\in L^0(\M)$, define the \emph{fraction commutator of $D$ and $Q$} by
\begin{equation}
 \Gamma[Q,D]:=sDb-aDt,
 \label{eq:fraction-commutator}
\end{equation}
with domain given by $\{\xi\in H: b\xi,t\xi\in\dom D\}$.
\end{definition}
The identities $ut=b$ and $su=a$ from \eqref{eq:fractions} give formally
\[
 s[D,u]t=sD(ut)-(su)D t=sD b-aD t=\Gamma[Q,D].
\]
That is, $s$ and $t$ remove the fractions in $u$ from the commutator expression.

Multiplication in $L^0(\M)\otimes_{\C}L^0(\M)^{\mathrm{op}}$ is given by
\[
 (u_1\otimes v_1^{\mathrm{op}})(u_2\otimes v_2^{\mathrm{op}})
 :=u_1u_2\otimes(v_2v_1)^{\mathrm{op}}.
\]
To state our first main result, we use the right multiplication maps in \cite[Section~3.1.1, equation~(3.1)]{CS05},
with the coefficient algebra enlarged from $\M\otimes_{\C}\M^{\mathrm{op}}$ to $L^0(\M)\otimes_{\C}L^0(\M)^{\mathrm{op}}$.

\begin{definition}\label{def:psi}
For $w\in L^0(\M)$, define the right multiplication map
\[
 \Psi_w:L^0(\M)\otimes_{\C}L^0(\M)^{\mathrm{op}}
 \longrightarrow L^0(\M)\otimes_{\C}L^0(\M)^{\mathrm{op}}
\]
by
\begin{equation}
 \Psi_w(k):=k(w\otimes1-1\otimes w^{\mathrm{op}})
 \label{eq:psi}
\end{equation}
\end{definition}

For $u,v\in L^0(\M)$, this gives
\[
 \Psi_w(u\otimes v^{\mathrm{op}})
 =uw\otimes v^{\mathrm{op}}-u\otimes(wv)^{\mathrm{op}}.
\]
For a tuple $X=(X_1,\ldots,X_d)\in\M^d$, we further write $\Psi_i:=\Psi_{X_i}$.

\Needspace{4\baselineskip}
\begin{theorem}\label{thm:delta}
Suppose $\Delta(X)=d$.
For $u\in L^0(\M)$ the following are equivalent:
\begin{enumerate}
\item $u\in\Div_{L^0(\M)}\A_X$.
\item\label{item:delta-tensors} There are $k_i\in L^0(\M)\otimes_{\C}L^0(\M)^{\mathrm{op}}$ such that
\begin{equation}
 \sum_{i=1}^d\Psi_i(k_i)=u\otimes1-1\otimes u^{\mathrm{op}}.
 \label{eq:tensor-criterion}
\end{equation}
\item\label{item:delta-all-fraction-tensors}
For every fraction pair $Q=(s,a,b,t)$ for $u$, there exist $\ell_1,\ldots,\ell_d\in L^0(\M)\otimes_{\C}L^0(\M)^{\mathrm{op}}$ such that
\begin{equation}\label{eq:all-fraction-tensor-criterion}
 \sum_{i=1}^d\Psi_i(\ell_i)
 =a\otimes t^{\mathrm{op}}-s\otimes b^{\mathrm{op}}.
\end{equation}
\item\label{item:delta-bounded-fraction-tensors}
There exist a fraction pair $Q=(s,a,b,t)$ for $u$ and tensors $\ell_1,\ldots,\ell_d\in\M\otimes_{\C}\M^{\mathrm{op}}$ satisfying \eqref{eq:all-fraction-tensor-criterion}.
\item\label{item:delta-finite-rank} There exist a fraction pair $Q=(s,a,b,t)$ for $u$ and a tuple $T=(T_1,\ldots,T_d)\in\F(H)^d$ such that
\begin{equation}
 \sum_{i=1}^d[\rho(X_i),T_i]=\Gamma[Q,P_\Omega].
 \label{eq:certificate}
\end{equation}
\end{enumerate}
In~\textup{(\ref{item:delta-tensors})}, $k_1,\ldots,k_d$ are uniquely determined by $u$.
The tensors $\ell_1,\ldots,\ell_d$ in \textup{(\ref{item:delta-all-fraction-tensors})} and \textup{(\ref{item:delta-bounded-fraction-tensors})} and the tuple $T$ in~\textup{(\ref{item:delta-finite-rank})} are uniquely determined by the fraction pair $Q$.
\end{theorem}

Here $P_\Omega$ denotes the orthogonal projection of $H$ onto $\C\Omega$.
By Definition~\ref{def:fraction-commutator},
\[
\Gamma[Q,P_\Omega]=sP_\Omega b-aP_\Omega t\in\F(H).
\]
The left-hand sides of equations~\eqref{eq:tensor-criterion} and~\eqref{eq:certificate} are values of degree-one Hochschild differentials with the corresponding coefficient bimodules; see Remark~\ref{rem:boundary-interpretation}.

\begin{remark}
When the equivalent assertions of Theorem~\ref{thm:delta} hold, the tensors in assertion~\textup{(\ref{item:delta-tensors})} are the free difference quotients $k_i=\partial_i u$ of $u$ with respect to $X_1,\ldots,X_d$;
see the implication $(1)\implies(2)$ in its proof.
The implication $(2)\implies(3)$ follows as $\Psi_i$ is a right multiplication map.
The implication $(3)\implies(4)$ is obtained by clearing denominators in $\ell_i$.
The implication $(4)\implies(5)$ is due to the identification of the tensors with finite-rank operators.
The key input for $(5)\implies(1)$ is that $\Div_{L^0(\M)}\A_X$ is a division ring, as guaranteed under $\Delta(X)=d$ by \cite[Theorem~1.1]{MSY23}.
\end{remark}

\subsection{Analytic criteria}
Condition~\textup{(\ref{item:delta-finite-rank})} in Theorem~\ref{thm:delta} requires a fraction pair $Q$ for which the commutator equation \eqref{eq:certificate} has the (unique) finite-rank solution $T\in\F(H)^d$.
The next theorem gives a rationality criterion that is valid for every fraction pair.
That is, \eqref{eq:certificate} admits a finite-rank solution either for every fraction pair or for none.

To formulate this analytic criterion, we write $\M\Omega:=\{x\Omega:x\in\M\}\subset H$.
We use the following variant of the (liberation) dual systems in \cite[Definition~13.1]{Voi99}.
\begin{definition}\label{def:liberation-dual-system}
Let $X=(X_1,\ldots,X_d)\in\M^d$.
A \emph{liberation dual system} for $X$ is a tuple $(L_1,\ldots,L_d)$ of linear operators on $H$ such that
\[
 \M\Omega\subseteq\dom L_i\cap\dom L_i^*
 \qquad(1\leq i\leq d),
\]
and
\[
 [X_j,L_i]\xi=\delta_{ij}[X_i,P_\Omega]\xi
 \qquad(\xi\in\M\Omega,\ 1\leq i,j\leq d).
\]
We call this system \emph{normalized} if
\begin{equation}\label{eq:liberation-normalization}
 \sum_{i=1}^dL_i\xi=P_\Omega\xi
 \qquad(\xi\in\M\Omega).
\end{equation}
\end{definition}
Every liberation dual system can be normalized, as shown in Lemma~\ref{lem:liberation-normalization}.

\Needspace{6\baselineskip}
\begin{theorem}\label{thm:intrinsic-analytic}
Let $X=(X_1,\ldots,X_d)\in\M^d$, and suppose that $X_i-\lambda$ is invertible in $L^0(\M)$ for every $1\leq i\leq d$ and every $\lambda\in\C$.
Let $(L_1,\ldots,L_d)$ be a normalized liberation dual system for $X$ in the sense of Definition~\ref{def:liberation-dual-system}.
Then, for every $u\in L^0(\M)$, the following are equivalent:
\begin{enumerate}
\item $u\in\Div_{L^0(\M)}\A_X$.
\item\label{item:intrinsic-analytic-existence-criterion}  There exist a fraction pair $Q$ for $u$ and $T_1,\ldots,T_d\in\F(H)$ such that
\begin{equation}
 \sum_{i=1}^d[\rho(X_i),T_i]=\Gamma[Q,P_\Omega].
 \label{eq:intrinsic-analytic-criterion}
\end{equation}
\item\label{item:intrinsic-analytic-criterion} For every fraction pair $Q$ for $u$, there exist $T_1,\ldots,T_d\in\F(H)$ satisfying \eqref{eq:intrinsic-analytic-criterion}.
\item\label{item:intrinsic-analytic-fraction-commutators} For every fraction pair $Q$ for $u$ and every $i$,
\[
 \overline{\Gamma[Q,L_i]}\in[\rho(X_i),\F(H)].
\]
\end{enumerate}
For each fixed fraction pair $Q$, the tuple $(T_1,\ldots,T_d)$ in assertions~\textup{(\ref{item:intrinsic-analytic-existence-criterion})} and~\textup{(\ref{item:intrinsic-analytic-criterion})} is unique.
\end{theorem}

Here $[\rho(X_i),\F(H)]:=\{[\rho(X_i),T]:T\in\F(H)\}$ and $\overline{\Gamma[Q,L_i]}$ denotes the closure of the closable operator $\Gamma[Q,L_i]$.
The new characterization in assertion~\textup{(\ref{item:intrinsic-analytic-fraction-commutators})} can be understood as follows.
For any fraction pair $Q$, assertion~\textup{(\ref{item:intrinsic-analytic-fraction-commutators})} provides $T_i\in\F(H)$ with $[\rho(X_i),T_i]=\overline{\Gamma[Q,L_i]}$.
Summing over $i$ and using the normalization in equation~\eqref{eq:liberation-normalization}, we obtain
\[
 \sum_i[\rho(X_i),T_i]
 =\sum_i\overline{\Gamma[Q,L_i]}
 =\Gamma[Q,P_\Omega],
\]
so $(T_1,\ldots,T_d)$ is the solution in assertion~\textup{(\ref{item:intrinsic-analytic-criterion})}.

Theorems~\ref{thm:delta} and~\ref{thm:intrinsic-analytic} do not require $X_1,\ldots,X_d$ to be normal.
For normal generators, we obtain two rationality criteria.
In Section~\ref{sec:conjugate-variable-criterion}, we apply Theorem~\ref{thm:intrinsic-analytic} to normal generators with nonatomic spectral distributions and $L^2$ liberation gradients.
The case of generators admitting conjugate variables follows as a corollary.

Our third main result further treats normal generators admitting a dual system and identifies the finite-rank solutions $T_i$ with the closures of the corresponding fraction commutators.
We use the following variant of the dual systems in \cite[Definition~5.1]{Voi98}.

\begin{definition}\label{def:dual-system}
Let $X=(X_1,\ldots,X_d)\in\M^d$.
A \emph{dual system} for $X$ is a tuple $(D_1,\ldots,D_d)$ of linear operators on $H$ such that
\begin{equation}\label{eq:analytic-hyp}
 \M\Omega\subseteq\dom D_i\cap\dom D_i^*,\qquad
 [X_j,D_i]|_{\M\Omega}=\delta_{ij}P_\Omega|_{\M\Omega}
 \qquad(1\leq i,j\leq d).
\end{equation}
\end{definition}

\Needspace{6\baselineskip}
\begin{theorem}\label{thm:analytic}
Suppose that $X_1,\ldots,X_d\in\M$ are normal operators generating $\M$, and let $(D_1,\ldots,D_d)$ be a dual system for $X$ in the sense of Definition~\ref{def:dual-system}.
For every $u\in L^0(\M)$ and every fraction pair $Q=(s,a,b,t)$ for $u$, the following are equivalent:
\begin{enumerate}
\item $u\in\Div_{L^0(\M)}\A_X$.
\item For every $i$, $\Gamma[Q,D_i]$ has finite rank.
\item For every $i$, $\overline{\Gamma[Q,D_i]}\in\F(H)$.
\end{enumerate}
When these equivalent conditions hold, the rank of each $\overline{\Gamma[Q,D_i]}$ is independent of $Q$.
\end{theorem}

Theorem~\ref{thm:analytic} identifies the finite-rank operators in Theorem~\ref{thm:intrinsic-analytic}\textup{(\ref{item:intrinsic-analytic-fraction-commutators})} with fraction-commutator closures for a dual system.
Lemma~\ref{lem:fundamental} constructs a normalized liberation dual system $(L_1,\ldots,L_d)$ from $(D_1,\ldots,D_d)$ and verifies regularity of every scalar shift.
For this system,
\[
 \overline{\Gamma[Q,L_i]}
 =[\rho(X_i),\overline{\Gamma[Q,D_i]}]
\]
whenever $\overline{\Gamma[Q,D_i]}$ is bounded.
Consequently, when these closures have finite rank, they are the finite-rank solution required in Theorem~\ref{thm:intrinsic-analytic}\textup{(\ref{item:intrinsic-analytic-existence-criterion})}, i.e.,
\[
 \sum_{i=1}^d[\rho(X_i),\overline{\Gamma[Q,D_i]}]=\Gamma[Q,P_\Omega].
\]
Conversely, normalization and Lemma~\ref{lem:finite-rank-identification} identify the finite-rank operators required in Theorem~\ref{thm:intrinsic-analytic}~\textup{(\ref{item:intrinsic-analytic-fraction-commutators})} with these closures.

A dual system for normal generators yields conjugate variables; the arcsine tuple of Remark~\ref{rem:liberation-without-conjugates} has a normalized liberation dual system but no conjugate variable, hence no dual system.
The arcsine tuple therefore satisfies the hypotheses of Theorem~\ref{thm:intrinsic-analytic} but not those of Theorem~\ref{thm:analytic}.
Hence the hypotheses of Theorem~\ref{thm:analytic} are strictly stronger than those of Theorem~\ref{thm:intrinsic-analytic}.

At present, we do not know whether normal generators admitting conjugate variables necessarily admit a dual system in the sense of Definition~\ref{def:dual-system}.
We present the two criteria in this order for convenience of exposition.
 
For normal operators admitting bounded dual systems, the fraction commutators are bounded and $\overline{\Gamma[Q,D_i]}=\Gamma[Q,D_i]$.
Theorem~\ref{thm:analytic} therefore recovers the criteria of Duchamp and Reutenauer and Linnell for free Haar unitaries, and of Miyagawa for free semicircular operators.

\subsection{Organization}\label{subsec:organization}

Section~\ref{sec:preliminaries} collects the algebraic preliminaries on division closures and the free field, followed by fractions of affiliated operators, finite-rank tensor realizations and the consequences of maximal $\Delta$.
Section~\ref{sec:maximal-delta} proves Theorem~\ref{thm:delta} using denominator clearing and the tensor realization of finite-rank operators.
Section~\ref{sec:liberation-dual-system-criterion} proves that liberation dual systems give maximal $\Delta$ when every $X_i-\lambda$ is regular, and establishes Theorem~\ref{thm:intrinsic-analytic} by recovering algebraic $L^2$ tensors.
Section~\ref{sec:dual-system-criterion} proves independence of the fraction pair, identifies finite-rank solutions with fraction-commutator closures, and constructs a normalized liberation dual system to deduce Theorem~\ref{thm:analytic} from Theorem~\ref{thm:intrinsic-analytic}\textup{(\ref{item:intrinsic-analytic-fraction-commutators})}.
Section~\ref{sec:conjugate-variable-criterion} constructs normalized liberation dual systems from $L^2$ liberation gradients and applies Theorem~\ref{thm:intrinsic-analytic} to normal generators with nonatomic spectral distributions.
It then obtains the conjugate-variable case as Corollary~\ref{prop:normal-conjugates} and gives an arcsine example separating the two hypotheses.
Section~\ref{sec:examples} uses a circular operator to illustrate the role of normality in Theorem~\ref{thm:analytic}.

% Source interfaces: MSY23-U002/U004; C_F0035; C_F0068 Sections 1--3.
\section{Preliminaries}\label{sec:preliminaries}

We first recall division closures, the free field and an algebraic tensor criterion for membership in a division subring.
We then discuss fractions of affiliated operators, realize algebraic tensors as finite-rank operators and recall the consequences of maximal $\Delta$.

\subsection{Division closures and the free field}\label{subsec:free-field}

\begin{definition}\label{def:rationality}
For a unital subalgebra $A$ of a unital algebra $B$, its \emph{division closure} $\Div_B A$ is the smallest unital subalgebra containing $A$ and closed under inversion of its elements that are invertible in $B$; see \cite[Exercise~7.1.4]{Coh85} and \cite[Definition~4.10]{MSY23}.
An affiliated operator $u$ is \emph{rational in $X$} if $u\in\Div_{L^0(\M)}\A_X$.
\end{definition}

The related notion of \emph{rational closure} is defined using entries of matrix inverses \cite[Definition~4.7]{MSY23}; see also \cite[Section~7]{Coh71}.
It contains the division closure and may be strictly larger in general \cite[Exercise~7.1.4]{Coh85}.
Since the division and rational closures coincide when $\Delta(X)=d$ (see Subsection~\ref{subsec:delta-tensor-maps}), we use \emph{rational in $X$} for membership in the division closure.

The free field $\freefield{x_1,\ldots,x_d}$ is the universal division ring of fractions of the free polynomial algebra $\C\langle x_1,\ldots,x_d\rangle$; see \cite[Section~4.1 and Theorem~4.1]{MSY23} and \cite[Corollary~7.5.14]{Coh06}.
It is the universal localization of this algebra at all full square matrices \cite[Theorem~7.5.13 and Corollary~7.5.14]{Coh06}.

The free difference quotients on noncommutative polynomials \cite[Section~3, pp.~192--193]{Voi98} extend uniquely to the free field \cite[Section~2.1, (4), pp.~298--299]{BD78}.
An explicit construction of these extensions on noncommutative rational functions is given in \cite[Section~4.4, Definition~4.4 and Corollary~4.11]{KV12}.
For $1\leq i\leq d$, we use the linear identification $u\otimes v\mapsto u\otimes v^{\mathrm{op}}$ to write the extended maps as
\[
 \partial_i:\freefield{x_1,\ldots,x_d}\longrightarrow
 \freefield{x_1,\ldots,x_d}\otimes_{\C}
 \freefield{x_1,\ldots,x_d}^{\mathrm{op}}.
\]
For $c\in\C$ and $u,v\in\freefield{x_1,\ldots,x_d}$, with $u\ne0$, they satisfy
\begin{equation}\label{eq:free-field-difference-rules}
\begin{aligned}
 \partial_i c&=0,& \partial_i x_j&=\delta_{ij}1\otimes1,\\
 \partial_i(uv)&=(1\otimes v^{\mathrm{op}})\partial_i u
                  +(u\otimes1)\partial_i v,&
 \partial_i(u^{-1})&=-(u^{-1}\otimes(u^{-1})^{\mathrm{op}})\partial_i u.
\end{aligned}
\end{equation}

\subsection{Coordinate maps over division subrings}

For a division subring $F$ in a unital algebra $U$, the proof below constructs, for $u\notin F$, a right $F$-linear coordinate map $\phi_u:U\longrightarrow F$ satisfying $\phi_u(1)=0$ and $\phi_u(u)=1$.
This map distinguishes $u\otimes_F1$ from $1\otimes_Fu$, which is recorded as the lemma below.
This is the key lemma used in the proof of Theorem~\ref{thm:delta} to force an affiliated operator to be a rational function in $X$.

\begin{lemma}\label{lem:equalizer}
Let $U$ be a nonzero unital algebra and let $F\subseteq U$ be a division subring containing the identity of $U$.
Then, for every $u\in U$,
\[
 u\otimes_F1=1\otimes_Fu\text{ in }U\otimes_FU
 \quad\Longleftrightarrow\quad u\in F.
\]
\end{lemma}
\begin{proof}
If $u\in F$, the identity follows from the balancing relation
\[
 (v_1a)\otimes_F v_2=v_1\otimes_F (av_2),\qquad v_1,v_2\in U,\quad a\in F.
\]
If $u\notin F$, the elements $1,u$ are right $F$-linearly independent: $v_1+uv_2=0$ with $v_2\ne0$ would give $u=-v_1v_2^{-1}\in F$.
Extend (1,u) to a basis of the right $F$-module $U$ by \cite[Theorem~IV.2.4]{Hun74}.
The coordinate map $\phi_u$ at $u$ is right $F$-linear (see \cite[Theorem~IV.2.1 \& Theorem~4.11]{Hun74}) and satisfies $\phi_u(1)=0$, $\phi_u(u)=1$.
The map $(v_1,v_2)\mapsto\phi_u(v_1)v_2$ is middle linear and hence induces an additive map \cite[Definition~IV.5.1 and Theorem~IV.5.2]{Hun74}
\[
 U\otimes_FU\longrightarrow U,\qquad v_1\otimes_Fv_2\longmapsto\phi_u(v_1)v_2
\]
which sends $u\otimes_F1-1\otimes_Fu$ to $1$, so this tensor cannot vanish.
\end{proof}

\subsection{Fractions of affiliated operators}\label{subsec:rational-affiliated}

We now turn to fraction representations in $L^0(\M)$.
Affiliation means commutation with each unitary in the commutant, including invariance of the domain; see \cite[Section 7.2]{AP18} and \cite[Section 5.1]{MSY23}.

\begin{definition}\label{def:regular}
An element $s\in\M$ is \emph{regular} if it is invertible in $L^0(\M)$.
\end{definition}

Regularity is equivalent to injectivity and dense range on $H$ \cite[Lemmas 5.1--5.2]{MSY23}.

For $u\in L^0(\M)$, a fraction pair is given by
\begin{equation}
 s:=(1+uu^*)^{-1/2},\quad a:=su,\qquad
 t:=(1+u^*u)^{-1/2},\quad b:=ut,
 \label{eq:canonical}
\end{equation}
where $s,a,b,t\in\M$, $s,t$ are regular, and $u=s^{-1}a=bt^{-1}$.
These facts follow from functional calculus and polar decomposition for affiliated operators \cite[Section 5.1]{MSY23}.

\begin{definition}\label{def:fraction-terminology}
For a fraction pair $Q=(s,a,b,t)$ of $u$ as in Definition~\ref{def:fraction-pair}, the representations $s^{-1}a$ and $bt^{-1}$ are called a \emph{left fraction} and a \emph{right fraction}, respectively; see \cite[Section~8.2.1]{Luc02}.

For $u_1,\ldots,u_m\in L^0(\M)$, a regular $s\in\M$ is a \emph{common left denominator} if $su_j\in\M$ for every $j$; a regular $t\in\M$ is a \emph{common right denominator} if $u_jt\in\M$ for every $j$; see \cite[Section~1]{Lin00} and \cite[Lemma~4.1]{Miy22}.
\end{definition}

Common denominators for finite families are used in \cite[Section~1, p.~2 of the arXiv version]{Lin00}.
The following proof iterates the two-element construction of \cite[Lemma~4.1]{Miy22}.

\begin{lemma}\label{lem:clearing}
For finite families $u_1,\ldots,u_m$ and $v_1,\ldots,v_n$ in $L^0(\M)$ there are regular $s,t\in\M$ such that
\[
 su_j\in\M\quad(1\leq j\leq m),\qquad
 v_kt\in\M\quad(1\leq k\leq n).
\]
\end{lemma}
\begin{proof}
Choose a regular $s\in\M$ with $su_1\in\M$ using \eqref{eq:canonical}.
Suppose $su_j\in\M$ for $1\leq j\leq\ell$, where $1\leq\ell<m$.
Using \eqref{eq:canonical}, write
\[
 u_{\ell+1}=s_1^{-1}a_1,\qquad ss_1^{-1}=s_2^{-1}a_2,
\]
where $s_1,s_2,a_1,a_2\in\M$ and $s_1,s_2$ are regular.
Then $s':=s_2s=a_2s_1$ is regular, and
\[
 s'u_j=s_2(su_j)\in\M\quad(1\leq j\leq\ell),\qquad
 s'u_{\ell+1}=a_2a_1\in\M.
\]
This completes the induction for the left denominators.
The proof for right denominators is similar.
\end{proof}

\begin{lemma}[{\cite[proof of Lemma~4.1]{Miy22}}]\label{lem:common}
If $s_1,s_2\in\M$ are regular, they have a common left multiple $xs_1=ys_2$ with $x,y\in\M$ regular.
If $t_1,t_2\in\M$ are regular, they have a common right multiple $t_1x'=t_2y'$ with $x',y'\in\M$ regular.
Both common multiples are regular.
\end{lemma}
\begin{proof}
Choose a left fraction $s_1s_2^{-1}=(s')^{-1}a$ using \eqref{eq:canonical}.
Thus we may take $x:=s'$ and $y:=a$.
Indeed, $s's_1=as_2$, and $a=s's_1s_2^{-1}$ is invertible in $L^0(\M)$, hence regular.
A right fraction $t_1^{-1}t_2=b_3t_3^{-1}$ gives the second identity, with $b_3=t_1^{-1}t_2t_3$ regular.
Products of regular elements are regular.
\end{proof}

Linnell \cite[Theorem~1.1]{Lin00} characterizes rational affiliated operators over a free group by finite-rank conditions on bounded operators obtained from left and right fractions.
Miyagawa \cite[Section~4]{Miy22} gives a semicircular analogue using the expressions
\[
 s r_i^*b-a r_i^*t=\Gamma[Q,r_i^*],
\]
where $Q=(s,a,b,t)$ is a fraction pair and $r_i^*$ are the bounded right annihilation operators.

\subsection{Algebraic tensors and finite-rank operators}\label{subsec:finite-rank-tensors}

Write $\F(H)$ for the finite-rank bounded operators on $H$.
For $x\in\M$, let $\rho(x)\in\B(H)$ denote the unique bounded extension of right multiplication
\[
 \rho(x)(y\Omega):=(yx)\Omega,\qquad y\in\M.
\]
We use the commutator convention $[S,T]:=ST-TS$.

Let $S_2(H)$ be the Schatten class of Hilbert--Schmidt operators, that is, the Hilbert space of operators $S\in\B(H)$ with $\operatorname{Tr}(S^*S)<\infty$, equipped with the Hilbert--Schmidt norm $\|S\|_{\operatorname{HS}}:=\operatorname{Tr}(S^*S)^{1/2}$.
Here $\operatorname{Tr}$ is the usual trace on $\B(H)$ \cite[Section 8.4]{AP18}.

Let $J$ be Tomita's conjugation operator on $H$, given by $J(x\Omega):=x^*\Omega$; see \cite[Section~7.1.1]{AP18} and \cite[p.~13]{MSY23}.
It satisfies $J^2=I$ and $\rho(x)=Jx^*J$.
Following \cite[Section~1.1, equation~(1.1)]{CS05}, define
\begin{equation}
 \Theta(u\otimes v)\xi:=u\langle Jv,\xi\rangle,
 \qquad \Theta:H\otimes_{\C}H\longrightarrow\F(H).
 \label{eq:theta}
\end{equation}
Here $u,v\in H=L^2(\M,\tau)\subset L^0(\M)$ and $\xi\in H$.
The map $\Theta$ is a linear bijection and extends to a unitary
\[
 \Theta:H\bar\otimes H\longrightarrow S_2(H).
\]
The actions in \cite[equation~(1.2)]{CS05} give, for $a,b\in\M$ and $u,v\in H$,
\[
 \Theta(au\otimes vb)=a\Theta(u\otimes v)b.
\]

Via the linear identification $u\otimes v\mapsto u\otimes v^{\mathrm{op}}$, we also regard $\Theta$ as a map on the corresponding subspace of $L^0(\M)\otimes_{\C}L^0(\M)^{\mathrm{op}}$.
For $u,v\in\M$, equation~\eqref{eq:theta} reads
\begin{equation}\label{eq:theta-bounded}
 \Theta(u\otimes v^{\mathrm{op}})=uP_\Omega v.
\end{equation}
For $x\in\M$ and $u,v\in H$, the right multiplication operator $\rho(x)$ satisfies
\begin{align*}
 \rho(x)\Theta(u\otimes v^{\mathrm{op}})
   &=\Theta(ux\otimes v^{\mathrm{op}}),\\
 \Theta(u\otimes v^{\mathrm{op}})\rho(x)
   &=\Theta(u\otimes(xv)^{\mathrm{op}}).
\end{align*}
Subtracting gives
\begin{equation}
 [\rho(x),\Theta(u\otimes v^{\mathrm{op}})]
   =\Theta(ux\otimes v^{\mathrm{op}}-u\otimes(xv)^{\mathrm{op}}).
 \label{eq:tensor-commutator}
\end{equation}

We extend $\tau$ to $H$ by $\tau(\xi):=\langle\Omega,\xi\rangle$.

\subsection{The invariant \texorpdfstring{$\Delta$}{Delta} and free-field realization}\label{subsec:delta-tensor-maps}

\Needspace{20\baselineskip}
\begin{definition}\label{def:delta}
The von Neumann tensor product $\M\bar\otimes\M^{\mathrm{op}}$, where $\M^{\mathrm{op}}$ is the opposite algebra, acts on $S_2(H)$ by
\[
 (a\otimes b^{\mathrm{op}})\cdot S:=aSb,\qquad a,b\in\M.
\]
For $X=(X_1,\ldots,X_d)\in\M^d$, set
\[
 \mathcal C_X:=
 \left\{T=(T_1,\ldots,T_d)\in\F(H)^d:
 \sum_{i=1}^d[\rho(X_i),T_i]=0\right\}.
\]
We call an element $T\in\mathcal C_X$ a \emph{finite-rank cycle}.
Such tuples represent Hochschild $1$-cycles with Hilbert--Schmidt coefficients; see \cite{CS05}.
Following \cite{CS05} and \cite[Section~3]{MSY23}, define
\begin{equation}
 \Delta(X):=d-\dim_{\M\bar\otimes\M^{\mathrm{op}}}
 \overline{\mathcal C_X}^{\operatorname{HS}}.
 \label{eq:delta-definition}
\end{equation}
The closure is taken in $S_2(H)^d$.
The dimension is the von Neumann dimension \cite[Definition~1.10]{Luc02} for the product trace $(a\otimes b^{\mathrm{op}})\mapsto\tau(a)\tau(b)$; see also \cite[Definition~2.10]{CS05}.
Thus $\Delta(X)\leq d$.
\end{definition}

We recall the finite-rank characterization of maximal $\Delta$ and the resulting free-field realization of the division closure, both from \cite{MSY23}.

\begin{remark}\label{rem:cycle-characterization}
By \cite[Section~3 and Remark~3.2]{MSY23}, for every $X\in\M^d$,
\[
 \Delta(X)=d
 \quad\Longleftrightarrow\quad \mathcal C_X=\{0\}
 \quad\Longleftrightarrow\quad
 \left\{T\in\F(H)^d:\sum_i[T_i,X_i]=0\right\}=\{0\}.
\]
\end{remark}

\begin{lemma}[{\cite[Theorems 1.1, 3.1 and 5.6]{MSY23}}]
\label{lem:free-field}
If $\Delta(X)=d$, evaluation at $X$ extends to an embedding of the free field $\freefield{x_1,\ldots,x_d}$ into $L^0(\M)$.
Its image is the division subring $\Div_{L^0(\M)}\A_X$.
\end{lemma}

If $\Delta(X)=d$, the rational closure and division closure of $\A_X$ in $L^0(\M)$ coincide \cite[Theorem~1.1 and Remark~5.10]{MSY23}.

% Complete proof of C_F0068 Part A, Sections 2--3; no Part B example.
\section{Rationality under maximal \texorpdfstring{$\Delta$}{Delta}}\label{sec:maximal-delta}

We prove Theorem~\ref{thm:delta} in this section.
By Definition~\ref{def:psi}, $\Psi_i$ is right multiplication by $X_i\otimes1-1\otimes X_i^{\mathrm{op}}$.
In particular,
\begin{equation}\label{eq:psi-left-linearity}
 \Psi_i(rk)=r\Psi_i(k),\qquad
 r,k\in L^0(\M)\otimes_{\C}L^0(\M)^{\mathrm{op}}.
\end{equation}

\begin{lemma}\label{lem:tensor-clearing}
Given $k_1,\ldots,k_d\in L^0(\M)\otimes_{\C}L^0(\M)^{\mathrm{op}}$, there are regular $s,t\in\M$ such that
\[
 (s\otimes t^{\mathrm{op}})k_i
 \in\M\otimes_{\C}\M^{\mathrm{op}}\qquad(1\leq i\leq d).
\]
\end{lemma}
\begin{proof}
For simplicity of notation, we give the calculation for simple tensors
\[
 k_i=u_i\otimes v_i^{\mathrm{op}},\qquad u_i,v_i\in L^0(\M).
\]
The same argument applies to general finite sums.
Applying Lemma~\ref{lem:clearing} to $u_1,\ldots,u_d$ gives a regular $s\in\M$ such that
\[
 su_i\in\M\qquad(1\leq i\leq d).
\]
Applying its right-denominator assertion to $v_1,\ldots,v_d$ gives a regular $t\in\M$ such that
\[
 v_it\in\M\qquad(1\leq i\leq d).
\]
Then
\[
 (s\otimes t^{\mathrm{op}})k_i
 =su_i\otimes(v_it)^{\mathrm{op}}
 \in\M\otimes_{\C}\M^{\mathrm{op}}.
\]
\end{proof}

\begin{lemma}\label{lem:tensor-injectivity}
Suppose $\Delta(X)=d$.
The map
\[
 (L^0(\M)\otimes_{\C}L^0(\M)^{\mathrm{op}})^d\longrightarrow
 L^0(\M)\otimes_{\C}L^0(\M)^{\mathrm{op}},\qquad
 (k_i)\longmapsto\sum_i\Psi_i(k_i)
\]
is injective.
\end{lemma}
\begin{proof}
Suppose $\sum_i\Psi_i(k_i)=0$.
Choose regular $s,t\in\M$ as in Lemma~\ref{lem:tensor-clearing} and set $T_i:=\Theta((s\otimes t^{\mathrm{op}})k_i)\in\F(H)$.
By \eqref{eq:tensor-commutator} and \eqref{eq:psi-left-linearity},
\[
 \sum_i[\rho(X_i),T_i]
 =\Theta\left((s\otimes t^{\mathrm{op}})\sum_i\Psi_i(k_i)\right)=0.
\]
Remark~\ref{rem:cycle-characterization} gives $T_i=0$ for every $i$.
Injectivity of $\Theta$ implies $(s\otimes t^{\mathrm{op}})k_i=0$.
Left multiplication by $s^{-1}\otimes(t^{-1})^{\mathrm{op}}$ then gives $k_i=0$.
\end{proof}

\begin{lemma}\label{lem:tensor-certificate}
Let $Q=(s,a,b,t)$ be a fraction pair for $u\in L^0(\M)$, and let $\ell_1,\ldots,\ell_d\in H\otimes_{\C}H$.
Set $T_i:=-\Theta(\ell_i)\in\F(H)$.
Then
\[
 \sum_i\Psi_i(\ell_i)
 =a\otimes t^{\mathrm{op}}-s\otimes b^{\mathrm{op}}
\]
if and only if
\[
 \sum_i[\rho(X_i),T_i]=\Gamma[Q,P_\Omega].
\]
\end{lemma}
\begin{proof}
By \eqref{eq:tensor-commutator} and \eqref{eq:theta-bounded},
\[
 \sum_i[\rho(X_i),T_i]
 =-\Theta\left(\sum_i\Psi_i(\ell_i)\right),
 \qquad
 \Gamma[Q,P_\Omega]
 =-\Theta(a\otimes t^{\mathrm{op}}-s\otimes b^{\mathrm{op}}).
\]
The equivalence follows from the injectivity of $\Theta$.
\end{proof}

By Lemma~\ref{lem:free-field}, $\Div_{L^0(\M)}\A_X$ is a division subring of $L^0(\M)$ whenever $\Delta(X)=d$.
We prove Theorem~\ref{thm:delta} through $(1)\implies(2)\implies(3)\implies(4)\implies(5)\implies(1)$.

\begin{proof}[Proof of Theorem~\ref{thm:delta}]
In this proof, set $U:=L^0(\M)$ and $F:=\Div_U\A_X$.

$(1)\implies(2)$.
By \cite[Theorem~1.1, (i) and (v)]{MSY23}, the assumption $\Delta(X)=d$ implies that $X_1,\ldots,X_d$ generate the free field: evaluation identifies $\freefield{x_1,\ldots,x_d}$ with $F$.
The free difference quotients from Subsection~\ref{subsec:free-field}, transported through this evaluation isomorphism, define linear maps
\[
 \partial_i:F\longrightarrow F\otimes_{\C}F^{\mathrm{op}}.
\]
They satisfy the rules in equation~\eqref{eq:free-field-difference-rules} with $x_j$ replaced by $X_j$.
Put $\delta v:=v\otimes1-1\otimes v^{\mathrm{op}}$.
Direct expansion gives
\[
 \delta(v_1v_2)=(1\otimes v_2^{\mathrm{op}})\delta v_1+(v_1\otimes1)\delta v_2,
 \qquad
 \delta(v^{-1})=-(v^{-1}\otimes(v^{-1})^{\mathrm{op}})\delta v.
\]
For constants, $\delta c=0$, while for generators, $\delta X_j=\Psi_j(1\otimes1)=\sum_i\Psi_i(\partial_iX_j)$.
We now show that this identity extends from constants and generators to all rational elements using the product and inverse rules.

If $\delta v_\ell=\sum_i\Psi_i(\partial_i v_\ell)$ for $\ell=1,2$, the product rules and equation~\eqref{eq:psi-left-linearity} give
\begin{align*}
 \delta(v_1v_2)
 &=(1\otimes v_2^{\mathrm{op}})\sum_i\Psi_i(\partial_i v_1)
   +(v_1\otimes1)\sum_i\Psi_i(\partial_i v_2)\\
 &=\sum_i\Psi_i\bigl((1\otimes v_2^{\mathrm{op}})\partial_i v_1
                  +(v_1\otimes1)\partial_i v_2\bigr)\\
 &=\sum_i\Psi_i\bigl(\partial_i(v_1v_2)\bigr).
\end{align*}
Similarly, if $0\ne v\in F$ satisfies $\delta v=\sum_i\Psi_i(\partial_i v)$, the inverse rules give
\begin{align*}
 \delta(v^{-1})
 &=-(v^{-1}\otimes(v^{-1})^{\mathrm{op}})\sum_i\Psi_i(\partial_i v)\\
 &=\sum_i\Psi_i\bigl(-(v^{-1}\otimes(v^{-1})^{\mathrm{op}})\partial_i v\bigr)\\
 &=\sum_i\Psi_i\bigl(\partial_i(v^{-1})\bigr).
\end{align*}
Together with linearity, this shows that the elements satisfying the identity form a division subalgebra of $F$ containing $\C$ and $X_1,\ldots,X_d$.
Since these generate $F$, the identity holds for every rational element $u\in F$.
This proves equation~\eqref{eq:tensor-criterion} with $k_i:=\partial_i u$; thus assertion~\textup{(\ref{item:delta-tensors})} holds.
Lemma~\ref{lem:tensor-injectivity} proves uniqueness of the tensors in assertion~\textup{(\ref{item:delta-tensors})}.

$(2)\implies(3)$.
Given tensors $k_i$ as in assertion~\textup{(2)}, for any fraction pair $Q=(s,a,b,t)$ set $\ell_i:=(s\otimes t^{\mathrm{op}})k_i$.
By equation~\eqref{eq:psi-left-linearity},
\[
 \sum_i\Psi_i(\ell_i)
 =(s\otimes t^{\mathrm{op}})
   (u\otimes1-1\otimes u^{\mathrm{op}})
 =a\otimes t^{\mathrm{op}}-s\otimes b^{\mathrm{op}},
\]
so assertion~\textup{(\ref{item:delta-all-fraction-tensors})} holds.

$(3)\implies(4)$.
Choose a fraction pair $Q=(s,a,b,t)$ for $u$ and corresponding tensors $\ell_i$ in assertion~\textup{(3)}.
By Lemma~\ref{lem:tensor-clearing}, there are regular $s',t'\in\M$ such that
\[
 (s'\otimes(t')^{\mathrm{op}})\ell_i
 \in\M\otimes_{\C}\M^{\mathrm{op}}.
\]
By equation~\eqref{eq:psi-left-linearity} and equation~\eqref{eq:all-fraction-tensor-criterion},
\[
 \sum_i\Psi_i\bigl((s'\otimes(t')^{\mathrm{op}})\ell_i\bigr)
 =s'a\otimes(tt')^{\mathrm{op}}-s's\otimes(bt')^{\mathrm{op}}.
\]
Since $(s's,s'a,bt',tt')$ is a fraction pair for $u$, assertion~\textup{(\ref{item:delta-bounded-fraction-tensors})} holds.

For fixed $Q$, uniqueness in assertions \textup{(\ref{item:delta-all-fraction-tensors})} and \textup{(\ref{item:delta-bounded-fraction-tensors})} follows from Lemma~\ref{lem:tensor-injectivity}, applied to the difference of two solutions of equation~\eqref{eq:all-fraction-tensor-criterion}.

$(4)\implies(5)$.
Let $Q$ and $\ell_i$ be as in assertion~\textup{(\ref{item:delta-bounded-fraction-tensors})}.
Lemma~\ref{lem:tensor-certificate} gives the finite-rank operators $T_i:=-\Theta(\ell_i)$ satisfying equation~\eqref{eq:certificate}.
Thus assertion~\textup{(\ref{item:delta-finite-rank})} holds.

$(5)\implies(1)$.
Suppose assertion~\textup{(\ref{item:delta-finite-rank})} holds.
Let $Q=(s,a,b,t)$ be the fraction pair and $T_1,\ldots,T_d\in\F(H)$ the operators satisfying equation~\eqref{eq:certificate}.
Via $\Theta$, define a linear map $\pi:\F(H)\longrightarrow U\otimes_FU$ by
\[
 \pi(\Theta(v_1\otimes v_2^{\mathrm{op}})):=v_1\otimes_Fv_2,
 \qquad v_1,v_2\in H.
\]
Note that $\Theta$ is a linear bijection, so $\pi$ is well defined.
For a rank-one operator $\Theta(v_1\otimes v_2^{\mathrm{op}})$, equation~\eqref{eq:tensor-commutator} gives
\begin{align*}
 \pi\bigl([\rho(X_i),\Theta(v_1\otimes v_2^{\mathrm{op}})]\bigr)
 &=(v_1X_i)\otimes_Fv_2-v_1\otimes_F(X_iv_2)\\
 &=0.
\end{align*}
The last equality uses $X_i\in F$ and the balancing relation
\[
 (uf)\otimes_Fw=u\otimes_F(fw),\qquad u,w\in U,\quad f\in F.
\]
Every finite-rank operator is a finite sum of rank-one operators, so linearity gives $\pi([\rho(X_i),T])=0$ for every $T\in\F(H)$.

Applying $\pi$ to equation~\eqref{eq:certificate} and using $\Theta(v_1\otimes v_2^{\mathrm{op}})=v_1P_\Omega v_2$ for $v_1,v_2\in\M$ (see equation~\eqref{eq:theta-bounded}), we obtain
\begin{align*}
 0&=\pi\left(\sum_i[\rho(X_i),T_i]\right)\\
  &=\pi(\Gamma[Q,P_\Omega])\\
  &=\pi(sP_\Omega b-aP_\Omega t)\\
  &=\pi\bigl(\Theta(s\otimes b^{\mathrm{op}})-\Theta(a\otimes t^{\mathrm{op}})\bigr)\\
  &=s\otimes_Fb-a\otimes_Ft.
\end{align*}
Hence $a\otimes_Ft=s\otimes_Fb$.

To clear the denominators, consider the map
\[
 \Phi_{s,t}:U\otimes_FU\longrightarrow U\otimes_FU,\qquad
 \Phi_{s,t}(u\otimes_Fw):=s^{-1}u\otimes_Fwt^{-1}.
\]
It is well defined because it preserves the balancing relation.
Applying $\Phi_{s,t}$ to $a\otimes_Ft=s\otimes_Fb$ yields
\[
 s^{-1}a\otimes_Ftt^{-1}=s^{-1}s\otimes_Fbt^{-1}.
\]
Since $u=s^{-1}a=bt^{-1}$, this is
\[
 u\otimes_F1=1\otimes_Fu.
\]
Lemma~\ref{lem:equalizer} therefore gives $u\in F=\Div_U\A_X$, so assertion~\textup{(1)} holds.

For fixed $Q$, let $(T_1,\ldots,T_d)$ and $(S_1,\ldots,S_d)$ be two finite-rank solutions of equation~\eqref{eq:certificate}.
Their right-hand sides are both $\Gamma[Q,P_\Omega]$, so subtraction gives
\begin{align*}
 \sum_i[\rho(X_i),T_i-S_i]
 &=\sum_i[\rho(X_i),T_i]-\sum_i[\rho(X_i),S_i]=0.
\end{align*}
Each $T_i-S_i$ has finite rank.
Since $\Delta(X)=d$, Remark~\ref{rem:cycle-characterization} implies $T_i-S_i=0$ for every $i$.
Hence the two tuples coincide, proving uniqueness of the finite-rank solution tuple for the fixed fraction pair $Q$.

\end{proof}

\begin{remark}\label{rem:boundary-interpretation}
Equip $L^0(\M)\otimes_{\C}L^0(\M)^{\mathrm{op}}$ with the $L^0(\M)$-bimodule structure defined, for $u,v\in L^0(\M)$ and $k\in L^0(\M)\otimes_{\C}L^0(\M)^{\mathrm{op}}$, by
\[
 u\cdot k\cdot v:=k(v\otimes u^{\mathrm{op}}).
\]
The degree-one Hochschild differential with these coefficients \cite[Section~9.1.1, p.~300]{Wei94} is the map
\[
 d_1:\bigl(L^0(\M)\otimes_{\C}L^0(\M)^{\mathrm{op}}\bigr)
 \otimes_{\C}L^0(\M)
 \longrightarrow L^0(\M)\otimes_{\C}L^0(\M)^{\mathrm{op}}
\]
given by
\[
 d_1(k\otimes u):=k\cdot u-u\cdot k
 =k(u\otimes1-1\otimes u^{\mathrm{op}}).
\]
Put $V_X:=\operatorname{span}_{\C}\{X_1,\ldots,X_d\}$.
Restricting the last tensor factor to $V_X$, we obtain
\[
 \partial_X:
 \bigl(L^0(\M)\otimes_{\C}L^0(\M)^{\mathrm{op}}\bigr)
 \otimes_{\C}V_X
 \longrightarrow L^0(\M)\otimes_{\C}L^0(\M)^{\mathrm{op}},
 \qquad \partial_X:=d_1|_{(L^0(\M)\otimes_{\C}L^0(\M)^{\mathrm{op}})\otimes_{\C}V_X}.
\]
By Definition~\ref{def:psi},
\[
 \partial_X\!\left(\sum_i k_i\otimes X_i\right)
 =\sum_i\Psi_i(k_i).
\]
This restriction to the linear span of the generators follows \cite[Section~3.1.1, equation~(3.1)]{CS05}, with $L^0(\M)\otimes_{\C}L^0(\M)^{\mathrm{op}}$ as the coefficient bimodule.
Thus condition~\textup{(\ref{item:delta-tensors})} in Theorem~\ref{thm:delta} requires $u\otimes1-1\otimes u^{\mathrm{op}}$ to lie in the image of $\partial_X$.

Similarly, equip $\F(H)$ with the $\A_X$-bimodule structure
\[
 x\cdot T\cdot y:=\rho(y)T\rho(x),
 \qquad x,y\in\A_X,\quad T\in\F(H).
\]
The degree-one Hochschild differential is then
\[
 d_1:\F(H)\otimes_{\C}\A_X\longrightarrow\F(H),
 \qquad d_1(T\otimes x):=T\cdot x-x\cdot T=[\rho(x),T].
\]
Since $\A_X$ is generated by $X_1,\ldots,X_d$, the identity
\[
 [\rho(xy),T]
 =[\rho(y),\rho(x)T]+[\rho(x),T\rho(y)]
\]
shows, by induction on monomials, that restricting the last tensor factor to $V_X$ gives the full image of $d_1$:
\[
 \operatorname{im}d_1
 =d_1\bigl(\F(H)\otimes_{\C}V_X\bigr)
 =\left\{\sum_{i=1}^d[\rho(X_i),T_i]:T_i\in\F(H)\right\}.
\]
Condition \eqref{eq:certificate} is therefore
\[
 d_1\!\left(\sum_iT_i\otimes X_i\right)=\Gamma[Q,P_\Omega].
\]
Thus $\Gamma[Q,P_\Omega]$ is a boundary, or equivalently its class vanishes in $H_0(\A_X,\F(H)):=\F(H)/\operatorname{im}d_1$.
\end{remark}

% Proof inputs: C_F0079 v4; C_F0068 Part A for maximal-Delta rationality.
\section{Liberation dual system criterion}
\label{sec:liberation-dual-system-criterion}

We prove Theorem~\ref{thm:intrinsic-analytic} in this section.

\begin{lemma}\label{lem:liberation-normalization}
If $X\in\M^d$ admits a liberation dual system, then it admits a normalized liberation dual system.
\end{lemma}
\begin{proof}
Let $(L_1,\ldots,L_d)$ be a liberation dual system for $X$.
Define operators $C,\widetilde L_1,\ldots,\widetilde L_d$, each with domain $\M\Omega$, by
\[
 C\xi:=\sum_{k=1}^dL_k\xi-P_\Omega\xi,
 \qquad \widetilde L_i\xi:=L_i\xi-\frac1dC\xi
 \qquad(\xi\in\M\Omega,\ 1\leq i\leq d).
\]
Since $X_j\M\Omega\subseteq\M\Omega$, the liberation relations give, on $\M\Omega$,
\[
 [X_j,C]
 =\sum_{k=1}^d[X_j,L_k]-[X_j,P_\Omega]
 =\sum_{k=1}^d\delta_{kj}[X_k,P_\Omega]-[X_j,P_\Omega]
 =0.
\]
Consequently, for $\xi\in\M\Omega$ and $1\leq i,j\leq d$,
\[
 [X_j,\widetilde L_i]\xi
 =[X_j,L_i]\xi=\delta_{ij}[X_i,P_\Omega]\xi,
 \qquad
 \sum_{i=1}^d\widetilde L_i\xi
 =\sum_{i=1}^dL_i\xi-C\xi
 =P_\Omega\xi.
\]
Since $\mathcal M\Omega\subseteq\operatorname{dom}L_i^*$, we have for $\eta,\xi\in\M\Omega$,
\[
 \langle\eta,\widetilde L_i\xi\rangle
 =\left\langle
 L_i^*\eta-\frac1d\sum_{k=1}^dL_k^*\eta+\frac1dP_\Omega\eta,
 \xi\right\rangle.
\]
By the definition of the adjoint, \(\eta\in\operatorname{dom}\widetilde L_i^*\), and hence \(\mathcal M\Omega\subseteq\operatorname{dom}\widetilde L_i^*\).
Thus \(\widetilde L_1,\ldots,\widetilde L_d\) satisfy all the defining conditions.
\end{proof}

\begin{lemma}[{see \cite[Problem~III.5.18]{Kat95}}]\label{lem:finite-closure}
Let \(T\) be a densely defined closable operator on a Hilbert space \(H\).
If \(T\) has finite rank, then \(\overline T\in\F(H)\), with the same range and rank as \(T\).
\end{lemma}

\begin{lemma}\label{lem:tensor-l2-recovery}
Let $x\in\M$ and suppose that $x-\lambda$ is regular for every $\lambda\in\C$.
The map $\Psi_x$ of Definition~\ref{def:psi} is injective.
If $k\in L^0(\M)\otimes_{\C}L^0(\M)^{\mathrm{op}}$ satisfies
\[
 \Psi_x(k)\in H\otimes_{\C}H,
\]
then $k\in H\otimes_{\C}H$.
\end{lemma}
\begin{proof}
If $T\in\F(H)$ satisfies $[T,x]=0$, then $\operatorname{ran}T$ is a finite-dimensional invariant subspace for $x$.
A nonzero such subspace would contain an eigenvector, contradicting regularity of $x-\lambda$.
Thus $T=0$, so Remark~\ref{rem:cycle-characterization} gives $\Delta(x)=1$.
Lemma~\ref{lem:tensor-injectivity}, applied to the one-element tuple $(x)$, now shows that $\Psi_x$ is injective.

To prove the remaining assertion, put $U:=L^0(\M)$ and define the right multiplication maps
\[
 \begin{aligned}
 R_x:U&\longrightarrow U,& R_x(u)&:=ux,\\
 R_{x^{\mathrm{op}}}:U^{\mathrm{op}}&\longrightarrow U^{\mathrm{op}},&
 R_{x^{\mathrm{op}}}(v^{\mathrm{op}})&:=(xv)^{\mathrm{op}}.
 \end{aligned}
\]
Thus $\Psi_x=R_x\otimes I-I\otimes R_{x^{\mathrm{op}}}$.
Neither right multiplication map has a nonzero finite-dimensional invariant subspace: such a subspace would contain an eigenvector, contradicting regularity of $x-\lambda$.
Since $x\in\M$, both maps preserve $H$, and their restrictions to $H$ have the same property.

For linear maps $A:V\to V$ and $B:W\to W$ on complex vector spaces, $A\otimes I-I\otimes B$ is injective on $V\otimes_{\C}W$ if $B$ has no nonzero finite-dimensional invariant subspace.
Indeed, write a nonzero kernel vector as $\sum_{r=1}^n v_r\otimes w_r$ with $n$ minimal, so both families are linearly independent.
Projecting the second factor modulo $\operatorname{span}\{w_1,\ldots,w_n\}$ shows that this span is $B$-invariant, a contradiction.
The same conclusion holds with $A$ in place of $B$.

The induced linear maps on the quotient spaces are
\[
 \begin{aligned}
 \widetilde R_x:U/H&\longrightarrow U/H,&
 \widetilde R_x(u+H)&:=ux+H,\\
 \widetilde R_{x^{\mathrm{op}}}:U^{\mathrm{op}}/H&\longrightarrow U^{\mathrm{op}}/H,&
 \widetilde R_{x^{\mathrm{op}}}(v^{\mathrm{op}}+H)&:=(xv)^{\mathrm{op}}+H.
 \end{aligned}
\]
Projecting the first factor of $k$ modulo $H$ gives a vector in the kernel of
\[
 \widetilde R_x\otimes I-I\otimes R_{x^{\mathrm{op}}}
 \quad\text{on }(U/H)\otimes_{\C}U^{\mathrm{op}}.
\]
This map is injective by the preceding observation, so exactness of tensoring over $\C$ gives $k\in H\otimes_{\C}U^{\mathrm{op}}$.
Projecting the second factor modulo $H$ now gives a vector in the kernel of
\[
 (R_x|_H)\otimes I-I\otimes\widetilde R_{x^{\mathrm{op}}}
 \quad\text{on }H\otimes_{\C}(U^{\mathrm{op}}/H).
\]
This map is injective because $R_x|_H$ has no nonzero finite-dimensional invariant subspace.
It follows that $k\in H\otimes_{\C}H$.
\end{proof}

For a liberation dual system $(L_1,\ldots,L_d)$, define linear maps
\[
 \Phi_i:\M\otimes_{\C}\M^{\mathrm{op}}
 \longrightarrow\Hom_{\C}(\M\Omega,H),\qquad
 \Phi_i(c\otimes d^{\mathrm{op}}):=cL_id|_{\M\Omega}.
\]
The liberation commutator relations give, for $z\in\M\otimes_{\C}\M^{\mathrm{op}}$ and $s,t\in\M$,
\begin{equation}\label{eq:liberation-tensor-actions}
 \Phi_i(\Psi_j(z))=\delta_{ij}\Theta(\Psi_i(z)),
 \qquad
 \Phi_i((s\otimes t^{\mathrm{op}})z)=s\Phi_i(z)t
 \quad\text{on }\M\Omega.
\end{equation}
Indeed, the first identity on $c\otimes d^{\mathrm{op}}$ is $c[X_j,L_i]d=\delta_{ij}c[X_i,P_\Omega]d$, and the second follows from multiplication in the opposite factor.

\begin{proposition}\label{prop:liberation-delta}
Let $X=(X_1,\ldots,X_d)\in\M^d$, and suppose that $X_i-\lambda$ is regular for every $1\leq i\leq d$ and every $\lambda\in\C$.
If $X$ admits a liberation dual system, then $\Delta(X)=d$.
\end{proposition}
\begin{proof}
Let $T_1,\ldots,T_d\in\F(H)$ satisfy $\sum_i[\rho(X_i),T_i]=0$, and put $k_i:=\Theta^{-1}(T_i)$.
Equation~\eqref{eq:tensor-commutator} gives $\sum_i\Psi_i(k_i)=0$.
By Lemma~\ref{lem:tensor-clearing}, choose regular $s,t\in\M$ such that $m_i:=(s\otimes t^{\mathrm{op}})k_i\in\M\otimes_{\C}\M^{\mathrm{op}}$ for every $i$.
Equation~\eqref{eq:psi-left-linearity} yields $\sum_i\Psi_i(m_i)=0$.
Applying $\Phi_i$ and equation~\eqref{eq:liberation-tensor-actions}, then using density of $\M\Omega$, we obtain
\[
 \Theta(\Psi_i(m_i))=0.
\]
Hence $\Psi_i(m_i)=0$.
The regularity assumption and Lemma~\ref{lem:tensor-l2-recovery} give $m_i=0$.
Cancellation of $s\otimes t^{\mathrm{op}}$ gives $k_i=0$, and therefore $T_i=0$.
Remark~\ref{rem:cycle-characterization} now gives $\Delta(X)=d$.
\end{proof}

\begin{lemma}\label{lem:tensor-images-l2}
Let $X=(X_1,\ldots,X_d)\in\M^d$, and suppose that $X_i-\lambda$ is regular for every $1\leq i\leq d$ and every $\lambda\in\C$.
Let $(L_1,\ldots,L_d)$ be a liberation dual system for $X$.
If $k_1,\ldots,k_d\in L^0(\M)\otimes_{\C}L^0(\M)^{\mathrm{op}}$ satisfy
\[
 h:=\sum_{i=1}^d\Psi_i(k_i)\in\M\otimes_{\C}\M^{\mathrm{op}},
\]
then, for every $i$,
\[
 k_i\in H\otimes_{\C}H,
 \qquad \overline{\Phi_i(h)}=\Theta(\Psi_i(k_i)).
\]
\end{lemma}
\begin{proof}
Fix $i$, write $h=\sum_r c_r\otimes d_r^{\mathrm{op}}$ with $c_r,d_r\in\M$, and put $S_i:=\Phi_i(h)$ on $\M\Omega$.
The adjoint-domain condition in Definition~\ref{def:liberation-dual-system} gives
\[
 \M\Omega\subseteq\dom S_i^*,\qquad
 S_i^*\eta=\sum_r d_r^*L_i^*(c_r^*\eta)
 \quad(\eta\in\M\Omega).
\]
Thus $S_i$ is closable.
By Lemma~\ref{lem:tensor-clearing}, choose regular $s,t\in\M$ such that
\[
 m_j:=(s\otimes t^{\mathrm{op}})k_j
 \in\M\otimes_{\C}\M^{\mathrm{op}}
 \qquad(1\leq j\leq d).
\]
Applying $\Phi_i$ to $\sum_j\Psi_j(m_j)=(s\otimes t^{\mathrm{op}})h$ and using equation~\eqref{eq:liberation-tensor-actions} gives
\begin{equation}\label{eq:liberation-cleared-operator}
 sS_it=\Theta(\Psi_i(m_i))\quad\text{on }\M\Omega.
\end{equation}
The right-hand side has finite rank and $s$ is injective, so the restriction of $S_i$ to $t\M\Omega$ has finite rank.
Since $t$ has dense range, $t\M\Omega$ is dense in $H$.
This restriction is closable, and Lemma~\ref{lem:finite-closure} makes its closure a bounded finite-rank operator on $H$.
The operator $\overline{S_i}$ extends that closure and therefore equals it.
Set $h_i:=\Theta^{-1}(\overline{S_i})\in H\otimes_{\C}H$.
Extending equation~\eqref{eq:liberation-cleared-operator} by boundedness and density, and using the outer bimodule actions of $\Theta$, we obtain
\[
 (s\otimes t^{\mathrm{op}})h_i
 =\Psi_i(m_i)
 =(s\otimes t^{\mathrm{op}})\Psi_i(k_i).
\]
Cancellation of the regular outside factors gives $\Psi_i(k_i)=h_i\in H\otimes_{\C}H$.
Since every $X_i-\lambda$ is regular, Lemma~\ref{lem:tensor-l2-recovery} gives $k_i\in H\otimes_{\C}H$.
\end{proof}

\begin{proof}[Proof of Theorem~\ref{thm:intrinsic-analytic}]
By Proposition~\ref{prop:liberation-delta}, $\Delta(X)=d$.
We prove $(1)\implies(4)\implies(3)\implies(2)\implies(1)$.

$(1)\implies(4)$.
Fix any fraction pair $Q=(s,a,b,t)$ for $u$.
By Theorem~\ref{thm:delta}\textup{(\ref{item:delta-all-fraction-tensors})}, there are tensors $\ell_i\in L^0(\M)\otimes_{\C}L^0(\M)^{\mathrm{op}}$ satisfying equation~\eqref{eq:all-fraction-tensor-criterion}.
Set
\[
 h:=a\otimes t^{\mathrm{op}}-s\otimes b^{\mathrm{op}}
   =\sum_i\Psi_i(\ell_i).
\]
Lemma~\ref{lem:tensor-images-l2} gives $\ell_i\in H\otimes_{\C}H$ and $\overline{\Phi_i(h)}=\Theta(\Psi_i(\ell_i))$.
By definition, $\Phi_i(h)=-\Gamma[Q,L_i]|_{\M\Omega}$.
The common adjoint-domain condition makes $\Gamma[Q,L_i]$ closable by Lemma~\ref{lem:fraction-domains}.
Its closure extends the bounded, everywhere-defined operator $-\overline{\Phi_i(h)}$ and therefore equals it.
With $T_i:=-\Theta(\ell_i)\in\F(H)$, equation~\eqref{eq:tensor-commutator} now gives
\[
 \overline{\Gamma[Q,L_i]}
 =-\Theta(\Psi_i(\ell_i))
 =[\rho(X_i),T_i].
\]
Thus assertion~\textup{(4)} holds.

$(4)\implies(3)$.
Fix any fraction pair $Q=(s,a,b,t)$ for $u$ and choose $T_i\in\F(H)$ such that
$[\rho(X_i),T_i]=\overline{\Gamma[Q,L_i]}$ for every $i$.
By equation~\eqref{eq:liberation-normalization}, on $\M\Omega$ we have
\[
 \begin{aligned}
  \sum_i[\rho(X_i),T_i]
  &=\sum_i\overline{\Gamma[Q,L_i]}\\
  &=s\Bigl(\sum_iL_i\Bigr)b-a\Bigl(\sum_iL_i\Bigr)t\\
  &=sP_\Omega b-aP_\Omega t
   =\Gamma[Q,P_\Omega].
 \end{aligned}
\]
Both sides are bounded, so density extends this identity to $H$, proving assertion~\textup{(3)}.

$(3)\implies(2)$.
Every affiliated operator admits a fraction pair, so this is immediate.

$(2)\implies(1)$.
This is the implication $\textup{(\ref{item:delta-finite-rank})}\implies\textup{(1)}$ of Theorem~\ref{thm:delta}.

For each fixed fraction pair, uniqueness of the tuple $(T_1,\ldots,T_d)$ follows from Theorem~\ref{thm:delta}.
\end{proof}

% Provenance: C_F0031 and C_F0068 Part C, finite-rank closure and domains.
\section{Dual system criterion}
\label{sec:dual-system-criterion}

We prove Theorem~\ref{thm:analytic} by applying the liberation dual system criterion of Theorem~\ref{thm:intrinsic-analytic}.
We first prove that the existence of a finite-rank closure and its rank are independent of the fraction pair, and identify finite-rank solutions of equation~\eqref{eq:certificate} with the closures of the corresponding fraction commutators.
We then construct a normalized liberation dual system directly from the dual operators and apply assertion~\textup{(\ref{item:intrinsic-analytic-fraction-commutators})} of Theorem~\ref{thm:intrinsic-analytic}.

The following lemma supplies the domain facts needed for the fraction commutators.
\begin{lemma}\label{lem:fraction-domains}
Let \(D_1,\ldots,D_d\) be densely defined linear operators on \(H\) such that
\[\M\Omega\subseteq\dom D_i\cap\dom D_i^*\]
for \(1\leq i\leq d\).
For every \(u\in L^0(\M)\), every fraction pair \(Q\) of \(u\), and every \(i\), the operator \(\Gamma[Q,D_i]\) is closable and densely defined.
\end{lemma}
\begin{proof}
Fix \(i\) and write \(Q=(s,a,b,t)\).
Since \(b,t\in\M\) preserve \(\M\Omega\), the domain of \(\Gamma[Q,D_i]\) contains \(\M\Omega\).
For \(\eta\in\M\Omega\) and \(\xi\in\dom\Gamma[Q,D_i]\), the hypothesis of the lemma gives
\begin{equation}\label{eq:defect-adjoint}
 \langle\eta,\Gamma[Q,D_i]\xi\rangle
 =\left\langle
 b^*D_i^*(s^*\eta)-t^*D_i^*(a^*\eta),\xi
 \right\rangle.
\end{equation}
Thus \(\M\Omega\subseteq\dom(\Gamma[Q,D_i])^*\), so \(\Gamma[Q,D_i]\) is closable.
\end{proof}

The next lemma allows us to pass from one fraction pair to every fraction pair before applying Theorem~\ref{thm:intrinsic-analytic}.

\begin{lemma}\label{lem:fraction-ranks}
Let $D$ be a densely defined linear operator on $H$ such that
$\M\Omega\subseteq\dom D\cap\dom D^*$.
For $u\in L^0(\M)$ and any two fraction pairs $Q,Q'$ for $u$,
\[
 \overline{\Gamma[Q,D]}\in\F(H)
 \quad\Longleftrightarrow\quad
 \overline{\Gamma[Q',D]}\in\F(H).
\]
When these conditions hold, the two closures have the same rank.
\end{lemma}
\begin{proof}
Write $Q=(s,a,b,t)$ and $Q'=(s',a',b',t')$.
Lemma~\ref{lem:common} gives regular $x,y,p,q\in\M$ such that $xs=ys'$ and $tp=t'q$.
Since $xa=ya'$ and $bp=b'q$, we have on $\M\Omega$
\begin{equation}\label{eq:fraction-common-comparison}
 x\Gamma[Q,D]p=y\Gamma[Q',D]q.
\end{equation}
Suppose that $\overline{\Gamma[Q,D]}\in\F(H)$.
The left-hand side of equation~\eqref{eq:fraction-common-comparison} has finite rank, and $y$ is injective.
Thus the restriction of $\Gamma[Q',D]$ to the dense subspace $q\M\Omega$ has finite rank.
It is closable by Lemma~\ref{lem:fraction-domains}, so Lemma~\ref{lem:finite-closure} makes its closure a bounded finite-rank operator on all of $H$.
The closure of $\Gamma[Q',D]$ extends this operator and hence equals it.
Interchanging $Q$ and $Q'$ proves the equivalence.

Both closures are now bounded, so equation~\eqref{eq:fraction-common-comparison} extends to
\[
 x\overline{\Gamma[Q,D]}p=y\overline{\Gamma[Q',D]}q
 \qquad\text{on }H.
\]
The left factors $x,y$ are injective, while $p,q$ have dense range.
These multiplications preserve the rank of a bounded finite-rank operator, proving the rank equality.
\end{proof}

We next identify any finite-rank solution of equation~\eqref{eq:certificate} with the fraction commutators of a dual system.

\begin{lemma}\label{lem:finite-rank-identification}
Let $X=(X_1,\ldots,X_d)\in\M^d$, and let $(D_1,\ldots,D_d)$ be a dual system for $X$.
Let $Q$ be a fraction pair for $u\in L^0(\M)$.
If $T_1,\ldots,T_d\in\F(H)$ satisfy
\[
 \sum_j[\rho(X_j),T_j]=\Gamma[Q,P_\Omega],
\]
then $\overline{\Gamma[Q,D_i]}=T_i$ for every $i$.
\end{lemma}
\begin{proof}
Write $Q=(s,a,b,t)$ and put $\ell_j:=-\Theta^{-1}(T_j)\in H\otimes_{\C}H$.
By Lemma~\ref{lem:tensor-certificate},
\[
 \sum_j\Psi_j(\ell_j)=a\otimes t^{\mathrm{op}}-s\otimes b^{\mathrm{op}}.
\]
Lemma~\ref{lem:tensor-clearing} gives regular $s',t'\in\M$ such that
$\ell_j':=(s'\otimes(t')^{\mathrm{op}})\ell_j\in\M\otimes_{\C}\M^{\mathrm{op}}$ for every $j$.
Equation~\eqref{eq:psi-left-linearity} gives
\[
 \sum_j\Psi_j(\ell_j')
 =s'a\otimes(tt')^{\mathrm{op}}-s's\otimes(bt')^{\mathrm{op}}.
\]
For fixed $i$ and $\xi\in\M\Omega$, apply to this equality the linear map
\[
 \M\otimes_{\C}\M^{\mathrm{op}}\longrightarrow H,\qquad
 c\otimes e^{\mathrm{op}}\longmapsto cD_i(e\xi).
\]
All inputs of $D_i$ lie in $\M\Omega$, and the dual relations give
\[
 cX_jD_i(e\xi)-cD_i(X_je\xi)=\delta_{ij}cP_\Omega e\xi.
\]
Using equation~\eqref{eq:theta-bounded}, we obtain
\[
 -s'\Gamma[Q,D_i]t'\xi
 =\Theta(\ell_i')\xi=-s'T_it'\xi.
\]
Since $s'$ is injective, $\Gamma[Q,D_i]$ agrees with $T_i$ on the dense subspace $t'\M\Omega$.
Lemma~\ref{lem:fraction-domains} gives closability of $\Gamma[Q,D_i]$.
Its closure therefore extends the bounded operator $T_i$, whose domain is all of $H$, and hence equals $T_i$.
\end{proof}

For \(N,S\in\B(H)\) with \(N\) normal, the classical Fuglede theorem \cite[Theorem~I, p.~35]{Fug50} gives \(SN=NS\implies SN^*=N^*S\).
The following analogue for bounded linear maps from \(\M\) to \(H\) will be used in Lemma~\ref{lem:module-commutation}.

\begin{lemma}[{see \cite[Proposition~1]{CS77}}]\label{lem:normal-intertwining}
Let \(x\in\M\) be normal, and let \(T:\M\to H\) be a bounded linear map.
If \(T(xz)=xT(z)\) for every \(z\in\M\), then \(T(x^*z)=x^*T(z)\) for every \(z\in\M\).
\end{lemma}

\begin{lemma}\label{lem:module-commutation}
Let \(T:\M\Omega\to H\) be a closable linear operator.
Then \(z\mapsto T(z\Omega)\) is bounded from \(\M\) to \(H\).
If \(x_1,\ldots,x_m\in\M\) are normal and \(T(x_j\xi)=x_jT\xi\) for every \(j\) and \(\xi\in\M\Omega\), then
\[
 T(a\xi)=aT\xi
 \qquad(a\in W^*(x_1,\ldots,x_m),\ \xi\in\M\Omega).
\]
\end{lemma}
\begin{proof}
Identify \(\M\Omega\) with \(\M\).
The graph of \(T:\M\to H\) is norm closed because the Hilbert-space closure \(\overline T\) is closed and agrees with \(T\) on \(\M\Omega\).
The closed graph theorem therefore gives the asserted bound.
Lemma~\ref{lem:normal-intertwining} gives \(T(x_j^*\xi)=x_j^*T\xi\).
Products and linear combinations give the same relation for the unital $*$-algebra generated by \(x_1,\ldots,x_m\).
For \(a\in W^*(x_1,\ldots,x_m)\), Kaplansky density \cite[Theorem~2.3.1]{AP18} gives a uniformly bounded net \(a_\lambda\) in this algebra converging strongly to \(a\).
For fixed \(\xi\in\M\Omega\),
\[
 a_\lambda\xi\longrightarrow a\xi,\qquad
 T(a_\lambda\xi)=a_\lambda T\xi\longrightarrow aT\xi.
\]
Closedness of \(\overline T\) and \(a\xi\in\M\Omega\) give \(T(a\xi)=aT\xi\).
\end{proof}

% Provenance: Yin_C_F0024-P1, Sections 2--6, and Yin_C_F0033-P1
% specialized to M Omega; the proof of the criterion uses C_F0024-P3
% (v2), identifying the C_F0063 certificates with actual fraction commutators.

We now construct a normalized liberation dual system from a dual system for normal generators and compare their fraction commutators.

\begin{lemma}\label{lem:fundamental}
Let $X_1,\ldots,X_d\in\M$ be normal operators generating $\M$, and let $(D_1,\ldots,D_d)$ be a dual system for $X$.
Then every $X_i-\lambda$, $\lambda\in\C$, is regular.
Define on $\M\Omega$
\begin{equation}\label{eq:dual-liberation-construction}
 K_i:=[\rho(X_i),D_i|_{\M\Omega}],\qquad
 C:=\sum_jK_j-P_\Omega,\qquad
 L_i:=K_i-\frac1d C.
\end{equation}
The tuple $(L_1,\ldots,L_d)$ is a normalized liberation dual system for $X$.
For every fraction pair $Q$ and every $i$, if $\overline{\Gamma[Q,D_i]}$ is bounded on $H$, then
\begin{equation}\label{eq:dual-liberation-fractions}
 \overline{\Gamma[Q,L_i]}
 =[\rho(X_i),\overline{\Gamma[Q,D_i]}].
\end{equation}
In particular, if $\overline{\Gamma[Q,D_i]}\in\F(H)$ for every $i$, then
\begin{equation}\label{eq:fundamental}
 \sum_{i=1}^d[\rho(X_i),\overline{\Gamma[Q,D_i]}]=\Gamma[Q,P_\Omega].
\end{equation}
\end{lemma}
\begin{proof}
For a spectral projection $e:=1_{\{\lambda\}}(X_i)$, the dual relation gives
$eP_\Omega e=e[X_i,D_i]e=0$ on $\M\Omega$.
Density and faithfulness of the trace give $e=0$.
Normality then implies that $X_i-\lambda$ is regular.

Since $\rho(X_i)^*=\rho(X_i^*)$ preserves $\M\Omega$, the adjoint-domain hypothesis gives
\[
 \M\Omega\subseteq\dom K_i^*,\qquad
 K_i^*\eta=D_i^*\rho(X_i^*)\eta-\rho(X_i^*)D_i^*\eta
 \quad(\eta\in\M\Omega).
\]
Consequently, $\M\Omega$ is contained in the adjoint domains of $C$ and each $L_i$ as well.
Left and right multipliers commute, and traciality gives
$[\rho(X_i),P_\Omega]=[X_i,P_\Omega]$ on $H$.
Jacobi's identity therefore gives on $\M\Omega$
\[
 [X_j,K_i]=[\rho(X_i),[X_j,D_i]]
 =\delta_{ij}[X_i,P_\Omega],\qquad [X_j,C]=0.
\]
It follows that $[X_j,L_i]=\delta_{ij}[X_i,P_\Omega]$ and $\sum_iL_i=P_\Omega$ on $\M\Omega$.
Thus $(L_1,\ldots,L_d)$ is a normalized liberation dual system.

The operator $C$ is closable, so Lemma~\ref{lem:module-commutation} gives
$C(a\xi)=aC\xi$ for all $a\in\M$ and $\xi\in\M\Omega$.
For $Q=(s,a,b,t)$, this implies
\[
 \Gamma[Q,C]\xi=(sb-at)C\xi=0\qquad(\xi\in\M\Omega).
\]
Right multipliers commute with the left multipliers in $Q$, and hence
\begin{equation}\label{eq:dual-liberation-on-domain}
 \Gamma[Q,L_i]\xi
 =\Gamma[Q,K_i]\xi
 =\rho(X_i)\Gamma[Q,D_i]\xi-\Gamma[Q,D_i]\rho(X_i)\xi
 \qquad(\xi\in\M\Omega).
\end{equation}
If $\overline{\Gamma[Q,D_i]}$ is bounded on $H$, the right-hand side extends to
$[\rho(X_i),\overline{\Gamma[Q,D_i]}]\in\B(H)$.
By Lemma~\ref{lem:fraction-domains}, $\Gamma[Q,L_i]$ is closable on its natural domain.
Its closure extends this bounded operator on all of $H$, proving equation~\eqref{eq:dual-liberation-fractions}.
Finally, summing equation~\eqref{eq:dual-liberation-fractions} on $\M\Omega$ and using $\sum_iL_i=P_\Omega$ gives equation~\eqref{eq:fundamental} by density.
\end{proof}

\begin{proof}[Proof of Theorem~\ref{thm:analytic}]
Let $(L_1,\ldots,L_d)$ be the normalized liberation dual system constructed in Lemma~\ref{lem:fundamental}.
That lemma also gives regularity of every scalar shift, so Theorem~\ref{thm:intrinsic-analytic} applies.
Fix $u\in L^0(\M)$ and a fraction pair $Q$ for $u$.

$\textup{(1)}\implies\textup{(3)}$.
By Theorem~\ref{thm:intrinsic-analytic}\textup{(\ref{item:intrinsic-analytic-fraction-commutators})}, there are $T_1,\ldots,T_d\in\F(H)$ such that
\[
 \overline{\Gamma[Q,L_i]}=[\rho(X_i),T_i]\qquad(1\leq i\leq d).
\]
Summing on $\M\Omega$ and using $\sum_iL_i=P_\Omega$, then extending by boundedness, gives
\[
 \sum_i[\rho(X_i),T_i]=\Gamma[Q,P_\Omega].
\]
Lemma~\ref{lem:finite-rank-identification} gives $\overline{\Gamma[Q,D_i]}=T_i$ for every $i$, proving assertion~\textup{(3)}.

$\textup{(3)}\implies\textup{(1)}$.
By Lemma~\ref{lem:fraction-ranks}, $\overline{\Gamma[Q',D_i]}\in\F(H)$ for every fraction pair $Q'$ of $u$ and every $i$.
Equation~\eqref{eq:dual-liberation-fractions} therefore gives
\[
 \overline{\Gamma[Q',L_i]}
 =[\rho(X_i),\overline{\Gamma[Q',D_i]}]
 \in[\rho(X_i),\F(H)].
\]
Thus assertion~\textup{(\ref{item:intrinsic-analytic-fraction-commutators})} of Theorem~\ref{thm:intrinsic-analytic} holds, and $u\in\Div_{L^0(\M)}\A_X$.

$\textup{(2)}\Longleftrightarrow\textup{(3)}$.
Lemma~\ref{lem:fraction-domains} gives closability, and Lemma~\ref{lem:finite-closure} gives the forward implication.
The reverse implication follows by restriction.
The rank independence is Lemma~\ref{lem:fraction-ranks}.
\end{proof}

% Proof inputs: C_F0082 and C_F0063 Appendix A for liberation contractions;
% C_F0025 for normal phase and C_F0080 for the arcsine obstruction.
\section{Applications to liberation gradients and conjugate variables}
\label{sec:conjugate-variable-criterion}

We first apply Theorem~\ref{thm:intrinsic-analytic} to normal generators with $L^2$ liberation gradients and nonatomic spectral distributions.
We then obtain the criterion for generators admitting conjugate variables as a corollary.

\subsection{Liberation gradients}\label{subsec:liberation-gradients}

We use the $L^2$ formulation of Voiculescu's liberation gradients in \cite[5.3--5.4 and Proposition~6.1]{Voi99}.
\begin{definition}\label{def:liberation-gradients}
Let $X=(X_1,\ldots,X_d)\in\M^d$ generate $\M$, and set
\[
 A_i:=W^*(X_i),\qquad B_i:=W^*(X_j:j\ne i),\qquad
 \mathcal P_i:=\operatorname{Alg}(A_i,B_i)
 \quad(1\leq i\leq d),
\]
where $B_1:=\C1$ if $d=1$.
We say that $X$ admits \emph{$L^2$ liberation gradients} if the derivations
\[
 \gamma_i:\mathcal P_i\subset H\longrightarrow
          \mathcal P_i\otimes_{\C}\mathcal P_i,\qquad
 \gamma_i(a):=a\otimes1-1\otimes a\quad(a\in A_i),\qquad
 \gamma_i(b):=0\quad(b\in B_i),
\]
are well-defined and $1\otimes1\in\dom\gamma_i^*$ for every $i$.
Here $\gamma_i^*$ is the Hilbert space adjoint for the dense domain $\mathcal P_i\Omega$ and the target $H\bar\otimes H$.
The liberation gradients $j_i:=j(A_i:B_i):=-\gamma_i^*(1\otimes1)\in H$ satisfy
\[
 j_i^*=-j_i,\qquad
 \tau(j_ip)=(\tau\otimes\tau)(\gamma_ip)
 \quad(p\in\mathcal P_i).
\]
\end{definition}
This condition is equivalent to finiteness of the liberation Fisher information
$\varphi^*(A_1;\ldots;A_d):=\sum_i\|j_i\|_2^2$ in \cite[Definition~9.1]{Voi99}.

The following construction supplies the operators needed for Theorem~\ref{thm:intrinsic-analytic}.
\begin{lemma}\label{lem:gradient-liberation}
If $X$ admits $L^2$ liberation gradients in the sense of Definition~\ref{def:liberation-gradients}, then $X$ admits a normalized liberation dual system.
\end{lemma}
\begin{proof}
Fix $i$, and define on $\mathcal P_i$
\[
 \delta_i:=\mathrm i\gamma_i,\qquad
 C_i^\circ(p\Omega):=(\id\otimes\tau)\delta_ip,\qquad
 \eta:=\mathrm i j_i.
\]
Then $\delta_i^*(1\otimes1)=\eta$.
The prescriptions on $A_i,B_i$ and the product rule show that $\delta_i$ is real:
$\delta_i(p^*)=(\delta_ip)^\dagger$, where $(u\otimes v)^\dagger:=v^*\otimes u^*$.
By \cite[Lemmas~4.6--4.7]{Mai15}, $\eta=\eta^*$, the algebraic tensor product $\mathcal P_i\otimes_{\C}\mathcal P_i$ is contained in $\dom\delta_i^*$, and
\[
 \langle C_i^\circ p,q\rangle=\langle p,q\eta-C_i^\circ q\rangle
 \qquad(p,q\in\mathcal P_i).
\]
In particular, both $\delta_i$ and $C_i^\circ$ are closable.
We adapt the argument in \cite[proof of Lemma~12]{Dab10}.
For $\delta_i p=\sum_r u_r\otimes v_r$, the prescriptions on $A_i,B_i$ and the product rule give
\[
 (\delta_i\otimes\id)\delta_i p-(\id\otimes\delta_i)\delta_i p
 =\mathrm i\sum_r u_r\otimes1\otimes v_r.
\]
Indeed, the mixed product-rule terms cancel, and the two sides agree on $A_i$ and $B_i$.
Since $C_i^\circ(\mathcal P_i)\subseteq\mathcal P_i$, applying $\id\otimes\tau\otimes\tau$ and using $\delta_i^*(1\otimes1)=\eta$ gives
\[
 (C_i^\circ)^2p=\sum_r u_r\tau(v_r\eta)+\mathrm iC_i^\circ p.
\]
The preceding integration by parts identity and the expansion of $\delta_i(p^*p)$ yield
\begin{equation}\label{eq:liberation-energy}
 \|C_i^\circ p\|_2^2
 =2\operatorname{Re}\langle C_i^\circ p,p\eta\rangle
  -\langle\delta_i(p^*p),1\otimes\eta\rangle
  -\langle\mathrm iC_i^\circ p,p\rangle.
\end{equation}

For $a=a^*\in\mathcal P_i$ with $\|a\|<1$, set $v:=a+\mathrm i(1-a^2)^{1/2}$.
The polynomial partial sums $v_n$ of its power series converge in operator norm.
The series for $\delta_i v_n$ and $\delta_i(v_n^*)$ converge in $H\bar\otimes H$, since $\|\delta_i(a^n)\|_2\leq n\|a\|^{n-1}\|\delta_i a\|_2$.
The product rule shows that $\delta_i(v_n^*v_n)$ converges, while $v_n^*v_n\to1$ in operator norm.
Closedness of $\overline{\delta_i}$ therefore gives $\delta_i(v_n^*v_n)\to0$.
Taking real parts and passing to the limit in \eqref{eq:liberation-energy}, with $\|v\|_2=1$, gives
\[
 \lim_n\|C_i^\circ v_n\|_2\leq2\|\eta\|_2+1.
\]
Apply the same argument to $v^*$; averaging via $a=(v+v^*)/2$, rescaling and splitting real and imaginary parts give, for some $K_i>0$,
\[
 \|C_i^\circ p\|_2\leq K_i\|p\|\qquad(p\in\mathcal P_i).
\]

The integration by parts identity shows that $\mathcal P_i\Omega\subseteq\dom(C_i^\circ)^*$ and
\[
 (C_i^\circ)^*(p\Omega)=p\eta-C_i^\circ(p\Omega),\qquad
 \|(C_i^\circ)^*(p\Omega)\|_2\leq(K_i+\|\eta\|_2)\|p\|
 \quad(p\in\mathcal P_i).
\]
Put $C_i:=\overline{C_i^\circ}$.
For $x\in\M$, Kaplansky density \cite[Theorem~2.3.1]{AP18} gives a uniformly bounded net $p_\lambda\in\mathcal P_i$ with $p_\lambda\Omega\to x\Omega$.
The estimate for $C_i^\circ p_\lambda$ gives a weakly convergent subnet in $H$.
The graph of $C_i$ is a closed linear subspace of $H\oplus H$ and hence is weakly closed, so $x\Omega\in\dom C_i$.
Applying the same argument to the bound for $(C_i^\circ)^*p_\lambda$ gives $x\Omega\in\dom C_i^*$.
Thus
\[
 \M\Omega\subseteq\dom C_i\cap\dom C_i^*.
\]
On $\mathcal P_i\Omega$, the product rule and the values of $\delta_i$ give
\[
 [C_i^\circ,X_j]=\mathrm i\delta_{ij}[X_i,P_\Omega].
\]
Since $\mathcal P_i\Omega$ is a core for $C_i$ and the right-hand side is bounded, approximation in the graph norm extends this identity to $\dom C_i$, with $X_j\dom C_i\subseteq\dom C_i$.
In particular it holds on $\M\Omega$.
Set $L_i:=\mathrm iC_i$.
Then both $\dom L_i$ and $\dom L_i^*$ contain $\M\Omega$, and
\[
 [X_j,L_i]|_{\M\Omega}=\delta_{ij}[X_i,P_\Omega]|_{\M\Omega}.
\]
This is a liberation dual system.
Lemma~\ref{lem:liberation-normalization} supplies a normalized system.
\end{proof}

For normal generators, nonatomicity supplies the remaining regularity hypothesis of Theorem~\ref{thm:intrinsic-analytic}.
\begin{proposition}\label{prop:liberation-gradients}
Let $X_1,\ldots,X_d\in\M$ be normal operators generating $\M$, each with nonatomic spectral distribution.
Suppose that $X$ admits $L^2$ liberation gradients in the sense of Definition~\ref{def:liberation-gradients}.
Then $\Delta(X)=d$, and $X$ admits a normalized liberation dual system.
For any such system, all the conclusions of Theorem~\ref{thm:intrinsic-analytic} hold.
In particular, for every $u\in L^0(\M)$ and every fraction pair $Q$ for $u$, we have $u\in\Div_{L^0(\M)}\A_X$ if and only if there exist $T_1,\ldots,T_d\in\F(H)$ such that
\[
 \sum_i[\rho(X_i),T_i]=\Gamma[Q,P_\Omega].
\]
For each fixed $Q$, the tuple $(T_1,\ldots,T_d)$ is unique whenever it exists.
\end{proposition}
\begin{proof}
Lemma~\ref{lem:gradient-liberation} gives a normalized liberation dual system.
Normality and nonatomicity imply that every $X_i-\lambda$, $\lambda\in\C$, is regular.
Theorem~\ref{thm:intrinsic-analytic} gives the criterion and uniqueness, and Proposition~\ref{prop:liberation-delta} gives $\Delta(X)=d$.
\end{proof}

\subsection{Conjugate variables}\label{subsec:ordinary-conjugates}

We now deduce the ordinary conjugate-variable case.
We use the following formulation, as in \cite[Proposition~3.5]{ACSY24}.

\begin{definition}\label{def:conjugate-variables}
Let $X=(X_1,\ldots,X_d)\in\M^d$ generate $\M$, and write
\[
 \A_X^\ast:=\C\langle X_1,\ldots,X_d,X_1^*,\ldots,X_d^*\rangle.
\]
We say that $X$ admits \emph{conjugate variables} if there exist non-commutative derivations
\[
 \partial_i:\A_X^\ast\subset H\longrightarrow H\bar\otimes H
 \qquad(1\leq i\leq d)
\]
such that
\begin{equation}
 \partial_i(X_j)=\delta_{ij}1\otimes1,\qquad
 1\otimes1\in\dom\partial_i^*
 \quad(1\leq i,j\leq d).
 \label{eq:normal-conjugate-condition}
\end{equation}
The vectors $\xi_i:=\partial_i^*(1\otimes1)\in H$ are called conjugate variables for $X$.
\end{definition}

The following commutator equation will identify the derivatives of the adjoint generators.
For a normal operator $x\in\M$ with nonatomic spectral distribution $\mu$, define in $L^2(\mu)\bar\otimes L^2(\mu)\subseteq H\bar\otimes H$
\begin{equation}\label{eq:normal-phase-kernel}
 v_x(z,w):=
 \begin{cases}
  (\bar z-\bar w)/(z-w),&z\ne w,\\
  0,&z=w,
 \end{cases}
 \qquad K_x:=\Theta(v_x)\in S_2(H).
\end{equation}

\begin{lemma}\label{lem:normal-adjoint-uniqueness}
Let $x\in\M$ be a normal operator with nonatomic spectral distribution $\mu$.
For every $c\in\C$, the equation
\[
 [K,x]=c[P_\Omega,x^*]
\]
has the unique solution $K=cK_x$ in $S_2(H)$.
\end{lemma}
\begin{proof}
The diagonal has $\mu\otimes\mu$ measure zero and $|v_x|=1$ off it.
Multiplication of its kernel gives
\[
 [K_x,x]=[P_\Omega,x^*].
\]
If two solutions differ by $S\in S_2(H)$, then $[S,x]=0$.
The classical Fuglede theorem gives $[S,x^*]=0$, so the positive compact operator $S^*S$ commutes with $x$.
If $S\ne0$, a positive eigenvalue of $S^*S$ has a nonzero finite-dimensional eigenspace invariant under $x$.
The restriction of $x$ to this space has an eigenvector, contrary to the absence of spectral atoms.
Thus $S=0$.
\end{proof}

We next verify the hypotheses of Proposition~\ref{prop:liberation-gradients}.
\begin{lemma}\label{lem:conjugate-liberation}
Let $X_1,\ldots,X_d\in\M$ be normal operators generating $\M$.
Suppose that $X$ admits conjugate variables in the sense of Definition~\ref{def:conjugate-variables}.
Then each $X_j$ has nonatomic spectral distribution, and $X$ admits $L^2$ liberation gradients.
\end{lemma}
\begin{proof}
Choose derivations $\partial_1,\ldots,\partial_d$ as in Definition~\ref{def:conjugate-variables}.
By \cite[Lemma~3.3]{ACSY24}, the algebraic tensor product $\A_X^\ast\otimes_{\C}\A_X^\ast$ is contained in $\dom\partial_i^*$, so each $\partial_i$ is closable.
Fix $i$, $\lambda\in\C$ and $\varepsilon>0$, and set
\[
 y:=X_i-\lambda,\qquad e:=1_{\{\lambda\}}(X_i),\qquad
 r_\varepsilon:=(y^*y+\varepsilon^2)^{-1},\qquad
 f_\varepsilon:=\varepsilon^2yr_\varepsilon.
\]
To check the domains, choose $R>\|y\|^2+\varepsilon^2$ and put $q:=1-R^{-1}(y^*y+\varepsilon^2)\in\A_X^\ast$.
Then $\|q\|<1$ and
\[
 r_\varepsilon=R^{-1}\sum_{n=0}^\infty q^n,\qquad
 \|\partial_i(q^n)\|_2
 \leq n\|q\|^{n-1}\|\partial_i q\|_2
 \quad(n\geq1).
\]
The partial sums converge in operator norm and their images under $\partial_i$ converge in $H\bar\otimes H$.
Thus $r_\varepsilon\in\dom\overline{\partial_i}$, with
\[
 \overline{\partial_i}r_\varepsilon
 =-r_\varepsilon\cdot
   \bigl(\partial_i(y^*)\cdot y+y^*\cdot\partial_i y\bigr)
   \cdot r_\varepsilon.
\]
Multiplying the partial sums by $\varepsilon^2y$ and using the product rule also gives $f_\varepsilon\in\dom\overline{\partial_i}$.
Since $y$ is normal, $yr_\varepsilon y^*=1-\varepsilon^2r_\varepsilon$, and hence
\begin{align*}
 \overline{\partial_i}f_\varepsilon
 &=\varepsilon^4r_\varepsilon\cdot\partial_i y\cdot r_\varepsilon
   -\varepsilon^2yr_\varepsilon\cdot\partial_i(y^*)\cdot yr_\varepsilon\\
 &=(\varepsilon^2r_\varepsilon)\otimes(\varepsilon^2r_\varepsilon)
   -(\varepsilon yr_\varepsilon)\cdot\partial_i(y^*)
     \cdot(\varepsilon yr_\varepsilon).
\end{align*}
Spectral calculus gives
\[
 \|\varepsilon^2r_\varepsilon\|\leq1,\qquad
 \|\varepsilon yr_\varepsilon\|\leq\tfrac12,\qquad
 \varepsilon^2r_\varepsilon\longrightarrow e,\qquad
 \varepsilon yr_\varepsilon\longrightarrow0
 \quad\text{strongly as }\varepsilon\downarrow0.
\]
Consequently, $\overline{\partial_i}f_\varepsilon\to e\otimes e$ in $H\bar\otimes H$.
On the other hand, $\|f_\varepsilon\|\leq\varepsilon/2$, so $f_\varepsilon\to0$ in $H$.
Closedness of $\overline{\partial_i}$ gives $e\otimes e=0$, and faithfulness of $\tau$ gives $e=0$.
Since $i$ and $\lambda$ were arbitrary, every $X_j$ has nonatomic spectral distribution.
Fix $i$, and write $\xi_i:=\partial_i^*(1\otimes1)$.
The map $\Psi_i$ acts on $H\bar\otimes H$ by the same formula.
On $\A_X^\ast$, define $\gamma_i^\circ:=\Psi_i\partial_i$.
Differentiating $X_jX_j^*=X_j^*X_j$ and applying Lemma~\ref{lem:normal-adjoint-uniqueness} gives $\partial_iX_j^*=\delta_{ij}v_{X_j}$, hence
\[
 \begin{aligned}
  \gamma_i^\circ(X_j)&=\delta_{ij}(X_j\otimes1-1\otimes X_j),\\
  \gamma_i^\circ(X_j^*)&=\delta_{ij}(X_j^*\otimes1-1\otimes X_j^*).
 \end{aligned}
\]
Since $\Psi_i$ commutes with the outer bimodule actions, $\gamma_i^\circ$ is a derivation with algebraic tensor values, and these identities make it anti-real.
The tensor-adjoint formula \cite[Lemma~4.3]{Mai15} gives
\[
 (\gamma_i^\circ)^*(1\otimes1)=[X_i^*,\xi_i].
\]
Consequently, $j_i:=[\xi_i,X_i^*]\in H$ is skewadjoint and satisfies
\[
 \tau(j_ip)=(\tau\otimes\tau)(\gamma_i^\circ p)
 \qquad(p\in\A_X^\ast).
\]

The corresponding word-moment formula holds for the algebras $\C[X_i,X_i^*]$ and $\C\langle X_j,X_j^*:j\ne i\rangle$.
By bounded Kaplansky approximation in each factor, it extends to $A_i=W^*(X_i)$ and $B_i=W^*(X_j:j\ne i)$, with the same $j_i$; see \cite[Remark~5.5(c)]{Voi99}.
To check well-definedness, let $r$ be a formal relation between $A_i$ and $B_i$ that evaluates to zero, and let $\widetilde\gamma_i(r)$ be its evaluated formal tensor derivative.
Applying the word-moment formula to $urv$ and using the product rule gives
\[
 (\tau\otimes\tau)\bigl(u\cdot\widetilde\gamma_i(r)\cdot v\bigr)=0
 \qquad(u,v\in\mathcal P_i).
\]
The other product-rule terms contain the vanishing factor $r$.
The faithful product-trace pairing therefore gives $\widetilde\gamma_i(r)=0$, so the derivation descends to $\mathcal P_i=\operatorname{Alg}(A_i,B_i)$.
The resulting $L^2$ moment identity gives
$1\otimes1\in\dom\gamma_i^*$ and $\gamma_i^*(1\otimes1)=-j_i$.
Thus $X$ admits $L^2$ liberation gradients in the sense of Definition~\ref{def:liberation-gradients}.
\end{proof}

\begin{corollary}\label{prop:normal-conjugates}
Let $X_1,\ldots,X_d\in\M$ be normal operators generating $\M$.
Suppose that $X$ admits conjugate variables in the sense of Definition~\ref{def:conjugate-variables}.
Then all the conclusions of Proposition~\ref{prop:liberation-gradients} hold.
In particular, $\Delta(X)=d$, and for any normalized liberation dual system for $X$, all the conclusions of Theorem~\ref{thm:intrinsic-analytic} hold.
\end{corollary}
\begin{proof}
Lemma~\ref{lem:conjugate-liberation} gives nonatomic spectral distributions and $L^2$ liberation gradients, so Proposition~\ref{prop:liberation-gradients} applies.
\end{proof}

% Proof input: C_F0080, monomial obstruction to the vacuum adjoint condition.
\begin{remark}\label{rem:liberation-without-conjugates}
The existence of a liberation dual system does not imply the existence of conjugate variables, even for normal generators with nonatomic spectral distributions.
Indeed, let
\[
 d\mu(x):=\frac{dx}{\pi\sqrt{1-x^2}}\quad(-1<x<1),\qquad
 \M:=L^\infty((-1,1),\mu),\qquad \tau(f):=\int f\,d\mu,
\]
and let $X(x):=x$.
Then $X=X^*$ generates $\M$ and has nonatomic spectral distribution, so every $X-\lambda$, $\lambda\in\C$, is regular.
Taking $L:=P_\Omega$ gives a normalized liberation dual system, since $[X,L]=[X,P_\Omega]$ and $L$ is bounded on $H$.
The liberation derivation relative to scalars is $a\mapsto a\otimes1-1\otimes a$, whose trace pairing is zero, so $j(W^*(X):\C)=0$ and Proposition~\ref{prop:liberation-gradients} applies.

Suppose that $X$ admits a conjugate variable in the sense of Definition~\ref{def:conjugate-variables}, and choose the corresponding derivation $\partial$.
Since $\partial X=1\otimes_{\C}1$, the product rule and the moments $\tau(X^{2r})=4^{-r}\binom{2r}{r}$ and $\tau(X^{2r+1})=0$ give
\[
 (\tau\otimes\tau)(\partial X^{2n+1})
 =4^{-n}\sum_{r=0}^{n}\binom{2r}{r}\binom{2n-2r}{n-r}=1,
\]
where the last equality follows by comparing coefficients in $(1-z)^{-1/2}(1-z)^{-1/2}=(1-z)^{-1}$.
On the other hand, $\|X^{2n+1}\|_2\to0$ by dominated convergence.
Thus $\xi:=\partial^*(1\otimes_{\C}1)\in H$ would satisfy
\[
 1=\bigl|(\tau\otimes\tau)(\partial X^{2n+1})\bigr|
 \leq\|\xi\|_2\,\|X^{2n+1}\|_2\longrightarrow0,
\]
a contradiction.
\end{remark}

\section{Examples}
\label{sec:examples}

A circular operator illustrates the distinction between the criteria in Theorems~\ref{thm:intrinsic-analytic} and~\ref{thm:analytic}.

\begin{remark}\label{rem:normality}
The normality assumption in Theorem~\ref{thm:analytic} cannot be omitted.
Let $l_j,r_j$, $j=1,2$, be the left and right creation operators on a full Fock space with the vacuum vector $\Omega$.
Their adjoints $l_j^*,r_j^*$ are the left and right annihilation operators, respectively.
The operators $S_j:=l_j+l_j^*$ form a standard free semicircular pair; see \cite[5.13]{Voi98}.
Set
\[
 c:=\frac{S_1+\mathrm iS_2}{\sqrt2},\qquad
 \M:=W^*(c),\qquad
 D:=-\frac{r_1^*-\mathrm i r_2^*}{\sqrt2}\in\B(H).
\]
The identities $[r_i^*,S_j]=\delta_{ij}P_\Omega$ give $[c,D]=P_\Omega$ and $[D,c^*]=0$.
Thus, for every fraction pair $Q=(s,a,b,t)$ for $c^*$,
\[
 \Gamma[Q,D]=s[D,c^*]t=0\in\F(H).
\]
However, $c^*\notin\Div_{L^0(\M)}\C[c]$: every element of this division closure commutes with $c$, whereas $[c,c^*]=\mathrm i[S_2,S_1]\ne0$.

The bounded dual $D$ also gives $\Delta(c)=1$ by \cite[Proposition~6.2]{MSY23}, so Theorem~\ref{thm:delta} still applies to $c$.
Since $c^*$ is not rational in $c$, none of its equivalent conditions holds for $c^*$.
In particular, the equation in Theorem~\ref{thm:delta}\textup{(\ref{item:delta-finite-rank})} has no finite-rank solution for any fraction pair for $c^*$.
For every fraction pair $Q=(s,a,b,t)$ for $c^*$, $\Gamma[Q,D]=0$ and thus fails to solve this equation:
\[
 [\rho(c),\Gamma[Q,D]]=0
 \ne s[P_\Omega,c^*]t=\Gamma[Q,P_\Omega].
\]
Here $[P_\Omega,c^*]\ne0$ since $c^*$ is centered and nonzero.
Regularity makes $s$ injective and $t$ have dense range, so $s[P_\Omega,c^*]t\ne0$.
\end{remark}

\bibliographystyle{amsalpha}
\bibliography{references}
\end{document}
